\documentclass[11pt]{article}

\setlength\textwidth{160mm}
\setlength\textheight{220mm}
\setlength\oddsidemargin{-3mm}
\setlength\headheight{0mm}
\setlength\topmargin{-10mm}

\usepackage{amsfonts,amsmath,amssymb,amsthm}
\usepackage[colorlinks=true,linkcolor=blue,citecolor=blue,urlcolor=blue]{hyperref}

\newtheorem{thm}{Theorem}[section]
\newtheorem{lem}[thm]{Lemma}
\newtheorem{prop}[thm]{Proposition}
\newtheorem{cor}[thm]{Corollary}
\newtheorem{fact}[thm]{Fact}
\theoremstyle{definition}
\newtheorem{defn}[thm]{Definition}
\newtheorem{example}[thm]{Example}
\theoremstyle{remark}
\newtheorem{rem}[thm]{Remark}

\newcommand{\R}{\mathbb{R}}
\newcommand{\Hc}{\mathcal{H}}
\newcommand{\Wc}{\mathcal{W}}
\newcommand{\hH}{\hat{H}}
\newcommand{\hb}{\hat{b}}
\newcommand{\hx}{\hat{x}}
\newcommand{\one}{\mathbf{1}}
\newcommand{\lmax}{\lambda_{\max}}
\newcommand{\lmin}{\lambda_{\min}}
\newcommand{\vt}{\vartheta(\bar G,w)}
\newcommand{\om}{\omega(G,w)}
\newcommand{\gr}[1]{\texttt{\detokenize{#1}}}
\DeclareMathOperator{\tr}{tr}
\DeclareMathOperator{\spn}{span}
\DeclareMathOperator{\aff}{aff}
\DeclareMathOperator{\rank}{rank}

\begin{document}

\title{Anchoring Cliques in Clique Wrappers\\[1ex]
{\large A supplement to ``A Least Squares Framework for the Maximum Weight
Clique Problem''}}

\author{Research notes prepared for Stanislav Busygin\\
with Claude (Anthropic)}

\date{September 26, 2026; revised October 10, 2026}

\maketitle

\begin{abstract}
The non-edge entries of a wrapping matrix are free parameters of the clique
wrapper. {\em Anchoring} spends them on a chosen maximal clique $Q$ so that
$Q$ satisfies the hypothesis of \cite[Theorem~2]{CW}; its indicator then
becomes a stationary point of the relaxed least squares program and lies on
the curve that the secular equation traces. We collect what was established
about anchoring, theoretically and in experiments with QUALEX-MS.
For a single clique we give the construction, show that
the indicator is then a point of the secular curve, and compute the curvature
of the wrapper towards every other clique, which explains why anchoring acts
as a local exchange around $Q$; the experiments confirm this and show that
anchoring pays off as a second pass on a good clique and not otherwise. For
several cliques the levels at which they can be anchored are forced: two
anchored cliques $Q_a,Q_b$ satisfy
$\mu_a W(Q_b\setminus Q_a)=\mu_b W(Q_a\setminus Q_b)$, so cliques of equal
weight share one level, and a shared level creates exactly the degenerate
case (multiple eigenvalues with vanishing linear form) that the secular
equation cannot see. We characterize when two cliques can be anchored
together. Over all wrappers the best value of the relaxed program is
$1/\vartheta(\bar G,w)$, the reciprocal of the Lov\'asz number, and the
standard wrapper gives Hoffman's bound on regular graphs; a clique indicator can be a global minimizer of the relaxed program only if
$\omega(G,w)=\vartheta(\bar G,w)$, in which case every Lov\'asz-optimal wrapper
anchors all maximum cliques at the top of the spectrum, in the hard case of
the trust region problem. Odd holes and odd antiholes admit no nontrivial
anchoring of all their maximum cliques, which with the strong perfect graph
theorem characterizes perfect graphs. For the quasigroup completion graphs
of Busygin and Pasechnik the standard wrapper is itself Lov\'asz-optimal on
every completable instance, and the completions are exactly the nonnegative
points of its degenerate sphere of global minimizers; this is why the hard
case cannot be exploited. Busygin's SAT01, which recasts many NP problems as
recognizing $\alpha<\vartheta=m$ in a weighted graph, generalizes this: an
equation wrapper is Lov\'asz-optimal exactly when $\vartheta=m$, which holds
exactly when the instance has a vector solution, and it is blind to the
contradictions outside the equations. For a proper Lov\'asz-optimal wrapper
the nonnegative points of its sphere of global minimizers are exactly the
maximum clique indicators, the vertices of largest norm of a polytope; the
hardness of SAT01 lies in finding them and not in a lack of properness,
since for exact cover the equation wrapper is proper and the nonnegative
points of its sphere are the exact covers. Every solution also lies on the
surface of the standard wrapper, which meets that sphere and the orthant
exactly in the solutions; the intersection of the sphere with a wrapper
surface is its intersection with a cone, explicit in two unit vectors, so
that of these conditions only nonnegativity is hard. Computations with a maximal clique
enumerator and with CSDP illustrate the results on the DIMACS graphs, on
random graphs, on random quasigroups with holes, and on SAT01 instances of
factoring and of the Hamilton cycle problem, where $\vartheta$ refutes the
smallest primes that propagation leaves open, with a vanishing margin, and
none of the non-Hamiltonian graphs.
\end{abstract}

\tableofcontents

\section{Introduction}

This note supplements \cite{CW}, whose notation it follows, and refers to
its statements as \cite[Theorem~2]{CW} and so on. The same notation is used in
QUALEX-MS \cite{B02}, whose trust region stage finds the stationary points
of \cite[Section~5]{CW}; Remark~\ref{rem:qualex} gives the dictionary.

\cite[Theorem~2]{CW} says that a maximal clique $Q$ is recognized through
a stationary point of the relaxed least squares program \cite[(7)]{CW} whenever
every vertex outside $Q$ is attached to $Q$ with the same weight, measured
through the wrapping matrix $H$ \cite[(6)]{CW}:
\begin{equation}
\label{eq:cw3}
\sum_{j \in Q} h_{ij} z_j = C z_i \quad \forall i \in V\setminus Q.
\end{equation}
For the standard wrapping matrix this is a strong combinatorial condition
(\cite[Corollary~1]{CW}), and it rarely holds. But the entries of $H$ on
non-adjacent pairs are free, and (\ref{eq:cw3}) is linear in them. {\em
Anchoring} $Q$ means choosing them so that (\ref{eq:cw3}) holds. (The
solver's sources and diagnostics call (\ref{eq:cw3}) ``Theorem 8''.)

The questions studied were:
\begin{enumerate}
\item What does anchoring one clique do to the stationary points, and is
it useful in QUALEX-MS? (Sections~\ref{sec:one} and~\ref{sec:exp}.)
\item Can several cliques, in particular all maximum weight cliques, be
anchored in one wrapper, and what does such a wrapper look like?
(Sections~\ref{sec:several} and~\ref{sec:all}.) For perfect graphs one
expects the Lov\'asz-optimal wrapper to anchor all maximum cliques at
$\mu=\lmax$; for general graphs the conditions look overdetermined.
\item How do these statements look on actual graphs?
(Section~\ref{sec:comp}.)
\item What do they say about SAT01 \cite{Bu07}, where deciding
$\alpha=\vartheta$ in a weighted graph captures many NP problems, and where
does its hardness lie once $\vartheta$ and a Lov\'asz-optimal wrapper are
known? (Sections~\ref{sec:sat01}, \ref{sec:hardness}, \ref{sec:surface}
and~\ref{sec:sat01comp}.)
\end{enumerate}

The main findings, in the order of the text:
\begin{itemize}
\item Anchoring one maximal clique is always possible, lowering entries
only (Proposition~\ref{prop:construct}); it makes the indicator of $Q$ the
point of the secular curve at the multiplier $\mu_Q=W(Q)-C$
(Lemma~\ref{lem:onsecular}). The curvature of the wrapper along the
direction from $x^Q$ to any other clique $Q'$ exceeds $\mu_Q$ by exactly
$\mu_Q (W(Q')-W(Q))/W(Q \triangle Q')$, whatever the free entries are
(Proposition~\ref{prop:curv}).
\item In QUALEX-MS, anchoring the clique found by a standard pass and
solving again never lost on 366 test graphs and improved 55 of 120 weighted
random graphs; anchoring the greedy clique lost far more often than it won
(Section~\ref{sec:exp}).
\item Two cliques anchored in one wrapper satisfy the {\em level identity}
$\mu_a W(Q_b\setminus Q_a)=\mu_b W(Q_a\setminus Q_b)$
(Theorem~\ref{thm:level}). Cliques of equal weight therefore share a level,
and a shared level forces an eigenvalue of $\hH$ whose eigenspace contains
the differences of their indicators and carries no linear form
(Theorem~\ref{thm:collision}): the degenerate case excluded in
\cite[Section~5]{CW}. Theorem~\ref{thm:two} characterizes when two
cliques can be anchored together; two maximum weight cliques always can
(Corollary~\ref{cor:twomax}), though some non-edges may have to be made
invisible (Proposition~\ref{prop:blind}).
\item $\max_H \min\{x^Tx : x\in\Wc^H\} = 1/\vt$
(Theorem~\ref{thm:theta}); for regular graphs the standard wrapper attains
Hoffman's ratio bound instead (Corollary~\ref{cor:hoffman}). A clique
indicator is a global minimizer of the relaxed program for some wrapper iff
$W(Q)=\om=\vt$ (Theorem~\ref{thm:global}); then every Lov\'asz-optimal
wrapper anchors every maximum clique at $\mu=\lmax(\hH)$, the hard case of
the trust region problem (Corollary~\ref{cor:hard}).
\item Anchoring all maximum cliques of an odd hole or an odd antihole
forces the level to $0$ (Theorem~\ref{thm:holes}), while graphs with
$\chi=\omega$ have an explicit Lov\'asz-optimal wrapper anchoring them all
(Proposition~\ref{prop:colour}). Hence a graph is perfect iff every induced
subgraph admits an anchoring of all its maximum cliques at a nonzero level
(Theorem~\ref{thm:perfect}).
\item Of the 31 DIMACS graphs with $n\leq1100$ and between 2 and 20000
maximum cliques, 11 do
not admit anchoring all of them at a nonzero level; of the remaining 20,
14 have $\vartheta=\omega$ and 6 have $\vartheta>\omega$
(Section~\ref{sec:comp}).
\item For the quasigroup completion graphs of \cite{BP03} the standard
wrapper is Lov\'asz-optimal on every completable instance, and the
completions are exactly the nonnegative points of its sphere of global
minimizers (Proposition~\ref{prop:qcp}); finding one is QCP. On random
quasigroups with holes its degenerate eigenspace was up to 55 times the
span of the completions (Section~\ref{sec:qwh}).
\item For a SAT01 instance with $m$ equations the equation wrapper $H_A$
generalizes that wrapper: it is Lov\'asz-optimal iff $\vartheta=m$, then
anchors every solution at the top, and its nonnegative global minimizers are
the exact covers of the equations, whatever the other contradictions
(Proposition~\ref{prop:sat01}). $\vartheta=m$ iff the instance has a vector
solution, and the Lov\'asz optima are the Gram matrices of the vector
solutions (Theorem~\ref{thm:vector}). After SAT01 propagation $\vartheta$
refutes four of the six primes below $2^{16}$ that propagation leaves open,
with a margin falling from $0.064$ to $7\cdot10^{-5}$, none of them without
the contradictions that propagation derives, and none of the non-Hamiltonian
graphs that survive it. Up to 15 bits the Lov\'asz optimum of a satisfiable
factoring instance is spanned by its solutions, a unique solution being the
only optimum, while the two solutions of $241\cdot251$ sit in an optimum of
rank 205 (Section~\ref{sec:sat01comp}).
\item For a proper Lov\'asz-optimal wrapper the nonnegative points of the
sphere of global minimizers are exactly the maximum clique indicators, the
vertices of largest norm of the polytope of nonnegative points of the top
eigenspace, and they exist iff $\omega=\vartheta$
(Proposition~\ref{prop:proper}). This does not make SAT01 easy: the
nonnegative points of the equation wrapper's sphere are the exact covers, so
deciding whether such a sphere has one is NP-complete, although $H_A$ is
proper when every contradiction lies in an equation and Lov\'asz-optimal on
every satisfiable instance (Corollary~\ref{cor:exactcover}). With
contradictions outside the equations a nonnegative non-clique point of the
optimal face makes every Lov\'asz-optimal wrapper improper
(Proposition~\ref{prop:improper}), and the Hamilton cycle instances have
such points (Section~\ref{sec:hardness}).
\item Every solution also lies on the surface $\Wc^{H_0}$ of the standard
wrapper, and the sphere of $H_A$, the orthant and $\Wc^{H_0}$ meet exactly in
the solutions (Proposition~\ref{prop:threesets}). The intersection of a
sphere with any wrapper surface is its intersection with a cone, given in
closed form by two unit vectors (Proposition~\ref{prop:cone}). Added to
QUALEX-MS's Douglas--Rachford stage, the surface solves one more factoring
instance with the orthant and adds nothing to the projection that respects
the contradictions outside the equations; without nonnegativity it finds
nothing at all (Section~\ref{sec:surface}).
\end{itemize}

\section{Setting}
\label{sec:setting}

\subsection{Wrapping matrices}

Let $G(V,E)$, $V=\{1,\ldots,n\}$, have positive vertex weights $w$, let
$z_i=\sqrt{w_i}$ and $W(S)=\sum_{i\in S} w_i$, and let $N(i)$ be the
neighborhood of $i$. For $S\subseteq V$ we write $\one_S$ for its
characteristic vector, $z_S$ for the vector with $(z_S)_i=z_i$ on $S$ and
$0$ elsewhere, and $x^S=z_S/W(S)$ for its indicator \cite[Definition~1]{CW}.

\begin{defn}
\label{def:family}
$\Hc(G,w)$ is the set of real symmetric $n\times n$ matrices $H$ with
$h_{ij}=z_i z_j$ whenever $i=j$ or $(i,j)\in E$. For $H\in\Hc(G,w)$ and a
non-adjacent pair $i\neq j$ we write
\[ h_{ij} = z_i z_j (1-\tau_{ij}), \]
and call $\tau_{ij}=\tau_{ji}$ the {\em defect} of the non-edge. The
standard wrapping matrix $H_0$ of \cite{CW} has $\tau\equiv 1$. The wrapping
matrices of \cite{CW}, which need $h_{ij}<z_iz_j$ off $E$, are the $H$ with
all defects positive; we call them {\em proper}.
\end{defn}

Every $H\in\Hc(G,w)$ defines the surface
$\Wc^H=\{x\in\R^n: x^THx=1,\ z^Tx=1\}$. The computation in the proof of
\cite[Proposition~2]{CW} involves only entries of $H$ inside the clique, so
{\em every clique indicator lies on $\Wc^H$ for every $H\in\Hc(G,w)$}.
Nonnegative defects are exactly what \cite[Proposition~1]{CW} needs.

The weighted Lov\'asz number of the complement of $G$ is
\begin{align}
\label{eq:thetaprimal}
\vt &= \min_{H\in\Hc(G,w)} \lmax(H) \\
\label{eq:thetadual}
    &= \max \{ \langle zz^T, X\rangle :\ X\succeq 0,\ \tr X = 1,\
       X_{ij}=0 \mbox{ for } i\neq j,\ (i,j)\notin E \},
\end{align}
where both optima are attained \cite{L79,GLS88,K94}. We call a matrix
attaining (\ref{eq:thetaprimal}) a {\em Lov\'asz-optimal} wrapping matrix
and a matrix attaining (\ref{eq:thetadual}) a {\em Lov\'asz optimum}.
Since the two programs have equal values, every Lov\'asz-optimal $H$ and
every Lov\'asz optimum $X$ satisfy the complementarity condition
\begin{equation}
\label{eq:compl}
(\vt\, I - H)\,X = 0,
\end{equation}
because $\tr((\vt I-H)X)=\vt-\langle zz^T,X\rangle=0$ ($H$ and $zz^T$
agree on the support of $X$) and both factors are positive semidefinite.

\begin{lem}
\label{lem:cliquevec}
For every $H\in\Hc(G,w)$ and every clique $Q$, $z_Q^THz_Q=W(Q)^2$. Hence
$\lmax(H)\geq W(Q)$ and $\om\leq\vt$; and $X_Q=z_Qz_Q^T/W(Q)$ is feasible
in (\ref{eq:thetadual}) with value $W(Q)$.
\end{lem}
\begin{proof}
$z_Q^THz_Q=\sum_{i,j\in Q} z_iz_jh_{ij}=\sum_{i,j\in Q}w_iw_j=W(Q)^2$ and
$\|z_Q\|^2=W(Q)$, so the Rayleigh quotient of $z_Q$ is $W(Q)$. $X_Q$ is
positive semidefinite with trace $1$ and is supported on $Q\times Q$, where
all pairs are adjacent; its value is $(z^Tz_Q)^2/W(Q)=W(Q)$.
\end{proof}

\subsection{The relaxed program and its stationary points}

For $H\in\Hc(G,w)$ let
\begin{equation}
\label{eq:v}
v(H) = \min \{ x^Tx :\ x\in\Wc^H \},
\end{equation}
the relaxed program \cite[(7)]{CW}. It is feasible (clique indicators) and its
objective is coercive, so the minimum is attained, and
$v(H)\leq 1/\om$.

As in \cite[Section~5]{CW} let $x^0=z/W(V)$, $P=I-zz^T/W(V)$, $\hH=PHP$,
$\hb=PHx^0$ and $s=1-(x^0)^THx^0$. Writing $x=x^0+\hx$ with $\hx\perp z$,
\[ x^Tx=\|x^0\|^2+\|\hx\|^2, \qquad
   x^THx=1 \iff q(\hx):=\hx^T\hH\hx+2\hb^T\hx=s. \]
In terms of the defects,
$s = W(V)^{-2}\sum w_iw_j\tau_{ij}$, the sum running over ordered
non-adjacent pairs $i\neq j$, so $s>0$ for proper wrapping matrices.
Throughout, eigenvalues of $\hH$ are those of its restriction to
$z^\perp$ (on $z$ itself $\hH$ vanishes).

With the multiplier convention of \cite{CW} (the multiplier of the
quadratic constraint written as $-1/\mu$), a point $x=x^0+\hx\in\Wc^H$ is
{\em stationary with multiplier $\mu\neq 0$} if
\begin{equation}
\label{eq:stat}
(\mu I-\hH)\,\hx=\hb \quad\mbox{on } z^\perp.
\end{equation}
In the eigenvector basis $\hH=U\,{\rm diag}(\lambda_i)\,U^T$, $c=U^T\hb$,
$y=U^T\hx$, (\ref{eq:stat}) reads $(\mu-\lambda_i)y_i=c_i$. For $\mu$ not
an eigenvalue, $x(\mu)=x^0+(\mu I-\hH)^{-1}\hb$ is the {\em secular curve},
and $x(\mu)\in\Wc^H$ iff $\mu$ solves the secular equation
\cite[(19)]{CW}.

\begin{lem}
\label{lem:global}
Let $x=x^0+\hx\in\Wc^H$.
\begin{enumerate}
\item[(a)] If $x$ is stationary with a multiplier $\mu>0$ such that
$\mu\geq\lmax(\hH)$, then $x$ is a global minimizer of (\ref{eq:v}).
\item[(b)] If $s>0$, $\lmax(\hH)>0$ and $x$ is a global minimizer of
(\ref{eq:v}), then $x$ is stationary with a multiplier
$\mu\geq\lmax(\hH)$.
\end{enumerate}
\end{lem}
\begin{proof}
(a) On $z^\perp$ the function $L(y)=\|y\|^2-\mu^{-1}(q(y)-s)$ has the
Hessian $2(I-\hH/\mu)\succeq 0$, so it is convex, and $\nabla L(\hx)=0$ by
(\ref{eq:stat}). Hence for every feasible $y$ (i.e. $q(y)=s$)
$\|y\|^2=L(y)\geq L(\hx)=\|\hx\|^2$.

(b) Let $r=\|\hx\|$. If some $y\in z^\perp$ with $\|y\|<r$ had
$q(y)\geq s$, then, since $q(0)=0<s$, $q(ty)=s$ for some $t\in(0,1]$, a
feasible point of norm $t\|y\|<r$. So $q<s$ on the open ball of radius $r$
and, by continuity, $\hx$ maximizes $q$ over the closed ball. By the
characterization of trust region solutions \cite{MS83} applied to $-q$
there is $\mu\geq0$ with $(\mu I-\hH)\hx=\hb$ and $\mu I-\hH\succeq 0$ on
$z^\perp$; thus $\mu\geq\lmax(\hH)>0$.
\end{proof}

\begin{rem}[QUALEX-MS]
\label{rem:qualex}
The trust region program of \cite{B02} (cf. \cite[Theorem~3]{CW}) maximizes
$x^TAx$, $A=H-w_{\min}I$, over $z^Tx=1$, $\|x\|=r$. On such a sphere
$x^Tx$ is constant, so its stationary points are those of $q$ on
$\|\hx\|=\sqrt{r^2-\|x^0\|^2}$; they satisfy (\ref{eq:stat}) with the same
$\hx$. The solver works with $\hat A=PAP=\hH-w_{\min}P$ and the multiplier
$\mu^*=\mu-w_{\min}$: the same eigenvectors and linear form, all eigenvalues
and multipliers shifted by $w_{\min}$. By Lemma~\ref{lem:global}, when $s>0$
and $\lmax(\hH)>0$ the global minimizers of (\ref{eq:v}) are the stationary
points with $\mu\geq\lmax(\hH)$, which are the global maximizers of the
trust region objective on their own sphere \cite{MS83}; in the solver's
terms, $\mu^*\geq\lmax(\hat A)$. Counts of ``eigenvalues above $\mu^*$'' reported by the solver are counts of
eigenvalues of $\hH$ above $\mu$.
\end{rem}

\subsection{The gauge}

\begin{lem}
\label{lem:gauge}
For $H\in\Hc(G,w)$ and $t\neq-1$ let $H_t=(1+t)H-t\,zz^T$. Then
$H_t\in\Hc(G,w)$ with defects $(1+t)\tau$, $\Wc^{H_t}=\Wc^H$ and
$v(H_t)=v(H)$; moreover $\hH_t=(1+t)\hH$, $\hb_t=(1+t)\hb$, $s_t=(1+t)s$,
and a point is stationary for $H$ with multiplier $\mu$ iff it is
stationary for $H_t$ with multiplier $(1+t)\mu$.
\end{lem}
\begin{proof}
On $E$ and the diagonal the entries of $H_t$ are
$(1+t)z_iz_j-tz_iz_j=z_iz_j$, and on a non-edge
$(1+t)z_iz_j(1-\tau_{ij})-tz_iz_j=z_iz_j(1-(1+t)\tau_{ij})$. For $z^Tx=1$,
$x^TH_tx=(1+t)x^THx-t$, which equals $1$ iff $x^THx=1$. Since $Pz=0$,
$\hH_t=(1+t)\hH$ and $\hb_t=(1+t)PHx^0-tPz(z^Tx^0)=(1+t)\hb$, and
$s_t=1-(1+t)(x^0)^THx^0+t=(1+t)s$. Multiplying (\ref{eq:stat}) by $1+t$
gives the last claim.
\end{proof}

So the wrapper surface depends only on the defects up to a common factor,
and the multipliers are defined up to the same factor. In particular the
uniform perturbation of the standard wrapper, $h_{ij}=-t z_iz_j$ on every
non-edge, is $(H_0)_t$ and changes nothing. For $1+t>0$
the position of a multiplier within the spectrum of $\hH$ and the ratio
\[ \rho(H)=\frac{\lmax(\hH)-\lmin(\hH)}{\|\hb\|} \]
are invariant; absolute spectral widths are not.

\section{Anchoring one clique}
\label{sec:one}

\begin{defn}
A maximal clique $Q$ is {\em anchored at level $\mu$} in $H\in\Hc(G,w)$ if
(\ref{eq:cw3}) holds with $C=W(Q)-\mu$.
\end{defn}

\begin{lem}
\label{lem:tau}
The following are equivalent:
\begin{enumerate}
\item[(a)] $Q$ is anchored at level $\mu$ in $H$;
\item[(b)] $\sum_{j\in Q\setminus N(i)} w_j\tau_{ij}=\mu$ for every
$i\in V\setminus Q$;
\item[(c)] $Hz_Q=(W(Q)-\mu)\,z+\mu\,z_Q$.
\end{enumerate}
\end{lem}
\begin{proof}
For $i\notin Q$,
$\sum_{j\in Q}h_{ij}z_j=z_i\sum_{j\in Q\cap N(i)}w_j
+z_i\sum_{j\in Q\setminus N(i)}w_j(1-\tau_{ij})
=z_i\bigl(W(Q)-\sum_{j\in Q\setminus N(i)}w_j\tau_{ij}\bigr)$,
which gives (a)$\iff$(b). For $i\in Q$, $(Hz_Q)_i=z_iW(Q)$, which is the
$i$-th entry of the right-hand side of (c) for every $\mu$; for $i\notin Q$
the $i$-th entry of (c) is (\ref{eq:cw3}).
\end{proof}

For the standard wrapper the sum in (b) is $W(Q\setminus N(i))$, the weight
of the non-neighbours of $i$ in $Q$, so $Q$ is anchored in $H_0$ iff that
weight is the same for all outside vertices, which is \cite[Corollary~1]{CW}.
Note that a clique that is not maximal can only be ``anchored'' at level
$0$: a vertex adjacent to all of $Q$ has an empty sum in (b).

\begin{lem}
\label{lem:onsecular}
Let $Q$ be a maximal clique and $\mu\neq0$. Then $Q$ is anchored at level
$\mu$ in $H$ iff $x^Q$ is a stationary point of (\ref{eq:v}) with
multiplier $\mu$. In that case $\hx^Q=x^Q-x^0$ satisfies
$(\mu I-\hH)\hx^Q=\hb$, and if $\mu$ is not an eigenvalue of $\hH$ then
$x^Q=x(\mu)$ is the point of the secular curve at $\mu$ and $\mu$ is a
root of the secular equation.
\end{lem}
\begin{proof}
Since $z^Tx^Q=1$, $Px^Q=x^Q-x^0=\hx^Q$, and $PHx^Q=PH(x^0+\hx^Q)=
\hb+\hH\hx^Q$. If $Q$ is anchored, Lemma~\ref{lem:tau}(c) gives
$Hx^Q=\frac{W(Q)-\mu}{W(Q)}z+\mu x^Q$, so $PHx^Q=\mu\hx^Q$, which is
(\ref{eq:stat}). Conversely, (\ref{eq:stat}) says $P(Hx^Q-\mu x^Q)=0$, i.e.
$Hz_Q=\mu z_Q+\kappa z$ for some $\kappa$; comparing an entry in $Q$ gives
$\kappa=W(Q)-\mu$, which is Lemma~\ref{lem:tau}(c).
\end{proof}

So anchoring is exactly what makes the secular equation of \cite{CW} find
$Q$: the curve $x(\mu)$ passes through $x^Q$ at $\mu=\mu_Q$. \cite[Theorem~2]{CW}
is the direction ``anchored $\Rightarrow$ stationary''; for proper wrappers
$\mu_Q>0$ by Lemma~\ref{lem:tau}(b).

\begin{prop}[the construction]
\label{prop:construct}
Let $Q$ be a maximal clique and $H\in\Hc(G,w)$. For $i\notin Q$ let
$C_i=z_i^{-1}\sum_{j\in Q}h_{ij}z_j$ and $m_i=W(Q\setminus N(i))>0$. For
any $C\leq\min_{i\notin Q}C_i$, the matrix $H'$ obtained by
\[ h'_{ij}=h'_{ji}=h_{ij}+\frac{C-C_i}{m_i}\,z_iz_j
\qquad (i\notin Q,\ j\in Q\setminus N(i)), \]
all other entries unchanged, lies in $\Hc(G,w)$, satisfies $H'\leq H$
entrywise, and anchors $Q$ at level $W(Q)-C$. If $H$ is proper so is $H'$.
\end{prop}
\begin{proof}
A non-edge between $i\notin Q$ and $j\in Q$ enters condition
(\ref{eq:cw3}) of $i$ only, so the conditions can be met one outside vertex
at a time; the correction raises $C_i$ by
$\frac{C-C_i}{m_i}\sum_{j\in Q\setminus N(i)}w_j=C-C_i$. The corrections
are nonpositive, which lowers entries and raises defects.
\end{proof}

QUALEX-MS implements this as \texttt{anchor\_wrapper()} with $C=\min_iC_i$,
optionally damping all corrections by a factor $\theta\in(0,1]$. For a single
clique every level can be reached; $C\leq\min_iC_i$ is what keeps the
corrections nonpositive.

\begin{prop}[curvature towards other cliques]
\label{prop:curv}
Let $Q$ be anchored at level $\mu$ in $H$, let $Q'\neq Q$ be any clique and
$u=x^{Q'}-x^Q$. Then $u\perp z$ and
\[ \frac{u^T\hH u}{u^Tu} = \mu\left(1+\frac{W(Q')-W(Q)}{W(Q\triangle Q')}
\right). \]
\end{prop}
\begin{proof}
Write $a=W(Q')$, $b=W(Q)$, $c=W(Q\cap Q')$. Since $z^Tu=0$,
$u^T\hH u=u^THu$. By Lemma~\ref{lem:cliquevec}, $z_{Q'}^THz_{Q'}=a^2$ and
$z_Q^THz_Q=b^2$; by Lemma~\ref{lem:tau}(c),
$z_{Q'}^THz_Q=(b-\mu)z_{Q'}^Tz+\mu z_{Q'}^Tz_Q=(b-\mu)a+\mu c$. Hence
\[ u^THu=\frac{a^2}{a^2}+\frac{b^2}{b^2}-\frac{2((b-\mu)a+\mu c)}{ab}
=\frac{2\mu(a-c)}{ab},\qquad
u^Tu=\frac1a+\frac1b-\frac{2c}{ab}=\frac{a+b-2c}{ab}, \]
and $a+b-2c=W(Q\triangle Q')$, so the quotient is
$2\mu(a-c)/(a+b-2c)=\mu(1+(a-b)/(a+b-2c))$.
\end{proof}

\begin{cor}
\label{cor:curv}
Let $Q$ be anchored at a level $\mu>0$.
\begin{enumerate}
\item[(a)] If some clique $Q'$ is heavier than $Q$, then $\lmax(\hH)>\mu$:
the anchored indicator is not a maximizer of the trust region program on
its sphere, and the excess curvature towards $Q'$ is the gain
$W(Q')-W(Q)$ per unit of exchanged weight $W(Q\triangle Q')$, times $\mu$.
\item[(b)] Towards a clique of the same weight the curvature equals $\mu$;
towards a lighter one it is smaller.
\end{enumerate}
\end{cor}

Part (a) is quantitative, not merely the observation that $x^{Q'}$ is a
feasible point of smaller norm: the directions of positive excess curvature
point to heavier cliques, most strongly to those that differ from $Q$ by a
small exchange with a gain. This is the local exchange behaviour measured in
Section~\ref{sec:exp}.

\section{Anchoring one clique in QUALEX-MS: experiments}
\label{sec:exp}

\subsection{Protocol and provenance}

A QUALEX-MS {\em pass} builds a wrapping matrix, runs the trust region
stage and refines the stationary points into cliques. The variants (named
as in \texttt{bench/variants}) were:
\begin{description}
\item[head] the standard pass with $H_0$;
\item[anc] ({\em cold}) one pass with $H_0$ anchored on the clique of the
greedy MIN heuristic; \textbf{anc05} the same with $\theta=0.5$;
\textbf{ancs} with the corrections dealt out to randomly permuted outside
vertices, so that (\ref{eq:cw3}) fails (a control); \textbf{ancp} re-anchored
on each improvement, at most 4 anchored passes;
\item[ancw] ({\em warm}) the standard pass, then one pass anchored on its
final clique; \textbf{ancws} its shuffled control; \textbf{ancwp} re-anchored
on each improvement, at most 4 anchored passes;
\item[icel] as ancw, then one line-searched step of $\lmax(H)$ minimization
along the gradient $2u_iu_j$ ($u$ the top eigenvector) on all non-edges,
entries clipped at $h_{ij}\leq0$, and re-anchoring (``fire and ice'');
\textbf{icels} with the step's entries shuffled over the non-edges;
\textbf{ices} a step on $\lmax(\hH)$ that only moves each outside vertex's
correction among its non-neighbours in $Q$, so that (\ref{eq:cw3}) stays
exact;
\item[lz] the standard pass replaced by a pass on the $\lmax$-stepped
standard wrapper; \textbf{lzw} that wrapper as a second pass after the
standard one; \textbf{lzws} its shuffled control.
\end{description}
The test sets were the 70 DIMACS graphs outside the ten largest; a set of
56 weighted graphs; uniform random graphs $G(n,p)$ with
$n\in\{200,300,500,1000\}$, $p\in\{0.25,0.5,0.7,0.8,0.9,0.95\}$, five of each,
unweighted (\texttt{ru}) and with integer weights uniform on $\{1..10\}$
(\texttt{rw}), for 70 of the weighted and 50 of the unweighted of which
cliquer \cite{NO03} proved the optimum; and for the largest runs the ten
largest DIMACS graphs and 24 random graphs with $n=2000$.

Results quoted from the experiments of 12--13 September 2026 were measured
before the benchmark tooling was put under version control; their logs,
and the generator of the 56-graph weighted set, were lost when the machine's
\texttt{/tmp} was wiped. Tables~\ref{tab:rw} and~\ref{tab:lz} are
reproducible: \texttt{bench/compare.py} computes them from the results
committed to the QUALEX-MS repository (directory \texttt{bench/runs},
solver built from commit 6cebb51), and they reproduce the corresponding
earlier figures exactly.

\subsection{The construction and the spectrum}

The construction is exact: on every anchored run the stationary point at
$\mu_Q$ returned the indicator to $10^{-15}$. The anchored indicator was
never the maximizer on its sphere: between 8 and 276 eigenvalues of $\hH$
lay above $\mu_Q$ (median 30 of 399 on DIMACS). Taking $C$ below
$\min_iC_i$ does not change this: each unit taken off $C$ raises $\mu_Q$ by
one but $\lmax$ by 1.2 to 2.6, and the count above $\mu_Q$ levels off at up
to $|Q|$ (the correction is supported on the $Q\times(V\setminus Q)$ block,
so it has at most $|Q|$ positive eigenvalues). By Theorem~\ref{thm:global}
below, no choice could have made the indicator a global maximizer on these
graphs unless $\vartheta=\omega$.

The corrections are large (smallest $t_i=(C-C_i)/m_i$: median $-7$ on DIMACS,
$-10$ on weighted graphs, down to $-132$), and forced: a vertex missing a
single vertex of $Q$ must put its whole correction on one entry. One
anchoring widens the absolute spectrum of $\hH$ (spectral radius $\times1.3$
to $2.1$ on unweighted graphs, $\times1.9$ to $3.6$ on weighted random ones,
9 graphs), but it inflates $\hb$ as well, and the gauge-invariant ratio
$\rho(H)$ decreased on 5 of the 9 graphs (\gr{brock200_1} 57 to 37,
\gr{g500_50_0} 88 to 40) and increased on 3 (\gr{p_hat300-3} 23 to 40).
Iterated anchoring does not compound: each pass rebuilds the standard
wrapper and anchors once, and $\lmax$ even fell as the anchor improved
(\gr{g2000_95_0} weighted: 127 standard, then 407, 266, 266, 261).

\subsection{Cold, warm and iterated anchoring}

{\em Cold} anchoring loses: DIMACS 2 better and 9 worse
(\gr{gen400_p0.9_65} $65\to52$), weighted random graphs 15 better and 67
worse. Neither $\theta=0.5$, nor a shuffled anchor (11/76), nor iterating
(31/50) rescued it. By Corollary~\ref{cor:curv}, anchoring makes the
neighbourhood of $Q$ attractive; around a poor clique that is a bad place.

{\em Warm} anchoring never lost on 366 graphs. It improved 3 DIMACS graphs
(\gr{keller5} $26\to27$, \gr{gen400_p0.9_55} $53\to55$, \gr{MANN_a45}
$342\to343$), 8 of the weighted-56, 15 of 120 unweighted and 55 of 120
weighted random graphs, at 1.5 to 1.9 times the time of one pass. A second
pass on a randomly perturbed wrapper of comparable cost improved 2, 0, 2 and
21 of the same.

{\em Iterating} it (ancwp) keeps climbing on dense weighted graphs: 35 of
the 120 weighted random graphs improved beyond the first anchored pass
(total weight $+179$, seven still improving at the cap of four passes),
3 of the weighted-56 (\gr{w_nrm_200_90_1}: $263\to267\to268\to270\to272$,
the optimum), 2 unweighted random graphs and no DIMACS graph; 38 of the 70
proven weighted optima were reached. On the largest graphs iterated warm
anchoring improved \gr{keller6} $53\to55$ and \gr{MANN_a81} $1096\to1098$,
3 of 12 unweighted $n=2000$ graphs (\gr{g2000_50_1} $15\to16$) and 8 of 12
weighted ones, with the largest gains seen (\gr{g2000_95_0} $891\to932$,
\gr{g2000_90_0} $530\to555$), at 2 to 4 times the time.

The gains are local: a median of 2 vertices dropped and 3 added, 95--97\% of
the anchor kept (91\% after iterating), against about 59\% for the gains of a
random wrapper. Whether meeting (\ref{eq:cw3}) exactly matters depends on
density: on weighted random graphs with $p\leq0.7$ the shuffled control did
as well (11 improved against 7; 37 proven optima each), but for $p\in\{0.9,
0.95\}$ only the exact anchor worked (38 against 6), and likewise on dense
unweighted graphs (12 against 1).

\subsection{Blending in one step of \texorpdfstring{$\lambda_{\max}$}{lambda-max} minimization}

Minimizing $\lmax(H)$ over $\Hc(G,w)$ is the Lov\'asz program
(\ref{eq:thetaprimal}); its optima are known to be highly degenerate, and
Sections~\ref{sec:several}--\ref{sec:all} explain why. One step of it was
tried instead of the full minimization. The top eigenvector $u$ of $H$ was
within about a percent of $z/\|z\|$, so the gradient $2u_iu_j$ has cosine
$0.98$ to $1.00$ with the gauge direction $z_iz_j$ of Lemma~\ref{lem:gauge},
which changes nothing. The remainder drives $\hb$ towards $0$: by
definition $\hb=PHz/W(V)$ vanishes iff $z$ is an eigenvector of $H$. The
ratio $\rho(H)$ grew by a median factor of $1.23$ to $1.65$ per suite, where
the shuffled control stayed at $1.00$.

Against one warm anchored pass the step (icel) raised the proven weighted
optima from 37 to 41 of 70, all four gains at $p\leq0.5$ (5 better and 0
worse there, against 0 and 1 for icels), and the unweighted ones from 43 to
44 of 50; on dense graphs it was 7 better and 9 worse, no different from its
control, and on DIMACS it gave back \gr{MANN_a45} $343\to342$. The
projected variant (ices) changed little (7 better, 5 worse on weighted random
graphs). Without anchoring (lz, lzw) the step is mostly the effect of any
large change of wrapper (Table~\ref{tab:lz}). The step also raised the
number of eigenvalues above $\mu_Q$ at the anchored point on 93 of the 120
weighted random graphs and lowered it on none (median 29.5 to 43): the
benefit of icel over ancw is not a more stable anchor.

\begin{table}[ht]
\begin{center}
\small
\begin{tabular}{|l|c|c|c|c|c|}
\hline
variant & proven optima & vs head & $p\leq0.5$ & $p=0.7$ & $p\geq0.8$ \\
        & (of 70) & better/worse & & & \\
\hline
head & 33 & -- & -- & -- & -- \\
ancw & 37 & 55/0 & 3/0 & 4/0 & 48/0 \\
icel & 41 & 59/0 & 7/0 & 4/0 & 48/0 \\
lz   & 35 & 33/36 & 6/6 & 4/9 & 23/21 \\
lzw  & 40 & 34/0 & 6/0 & 5/0 & 23/0 \\
lzws & 35 & 28/0 & 3/0 & 3/0 & 22/0 \\
\hline
icel vs ancw & & 12/9 & 5/0 & 0/0 & 7/9 \\
lzw vs lzws  & & 20/18 & 5/1 & 3/3 & 12/14 \\
\hline
\end{tabular}
\end{center}
\caption{\label{tab:rw} Weighted random graphs (120, of which 70 with proven
optima). Time relative to head: ancw 1.88, icel 2.20, lz 1.34, lzw 1.98,
lzws 1.76.}
\end{table}

\begin{table}[ht]
\begin{center}
\small
\begin{tabular}{|l|c|c|c|}
\hline
variant & DIMACS (70) & unweighted random (120) & weighted random (120) \\
\hline
lz   & 2/0 & 10/3 & 33/36 \\
lzw  & 1/0 & 11/0 & 34/0 \\
lzws & 3/0 & 7/0  & 28/0 \\
\hline
\end{tabular}
\end{center}
\caption{\label{tab:lz} One $\lmax$ step without anchoring, better/worse
against head. The DIMACS changes were \gr{DSJC1000.5} $14\to15$ (lz, lzw),
\gr{gen400_p0.9_55} $53\to54$ (lz) and $\to55$ (lzws), \gr{keller5}
$26\to27$ and \gr{MANN_a45} $342\to343$ (lzws). On the unweighted random
graphs all three reach 43 of the 50 proven optima (head 42). For
comparison, as second passes: a random wrapper 2/2/21, warm anchoring
3/15/55.}
\end{table}

\section{Anchoring several cliques}
\label{sec:several}

We say that $H$ anchors the maximal cliques $Q_1,\ldots,Q_k$ at levels
$\mu_1,\ldots,\mu_k$ if it anchors each $Q_a$ at $\mu_a$. By
Lemma~\ref{lem:tau} this is a system of linear equations in the defects,
one for each pair $(Q_a,i)$ with $i\notin Q_a$:
\begin{equation}
\label{eq:system}
\sum_{j\in Q_a\setminus N(i)} w_j\tau_{ij}=\mu_a .
\end{equation}
A non-edge $\{i,j\}$ can now occur in several equations, and the system can
be inconsistent. For $Q_a\neq Q_b$ we write $D_a=Q_a\setminus Q_b$,
$D_b=Q_b\setminus Q_a$ and $K=Q_a\cap Q_b$; for maximal cliques $D_a$ and
$D_b$ are nonempty.

\begin{thm}[level identity]
\label{thm:level}
If $H$ anchors two distinct maximal cliques $Q_a,Q_b$ at levels
$\mu_a,\mu_b$, then
\[ \mu_a\,W(Q_b\setminus Q_a) = \mu_b\,W(Q_a\setminus Q_b). \]
\end{thm}
\begin{proof}
A vertex $i\in D_b$ lies outside $Q_a$ and is adjacent to all of $K$, so
its non-neighbours in $Q_a$ lie in $D_a$ and (\ref{eq:system}) reads
$\sum_{j\in D_a\setminus N(i)}w_j\tau_{ij}=\mu_a$. Multiplying by $w_i$ and
summing over $i\in D_b$ gives
$\sum w_iw_j\tau_{ij}=\mu_aW(D_b)$, the sum running over the non-edges
between $D_a$ and $D_b$. The same computation from the side of $Q_b$ gives
the same sum equal to $\mu_bW(D_a)$.
\end{proof}

A second proof: apply Proposition~\ref{prop:curv} to $u=x^{Q_b}-x^{Q_a}$
from both ends. The Rayleigh quotient of $u$ is
$\mu_a(1+(W_b-W_a)/W(\triangle))=\mu_b(1+(W_a-W_b)/W(\triangle))$, with
$W_a=W(Q_a)$ and $W(\triangle)=W(D_a)+W(D_b)$, which rearranges to the
identity.

\begin{cor}
\label{cor:equal}
Cliques of equal weight anchored in one wrapper share one level; in
particular all maximum weight cliques do. Pairwise disjoint anchored cliques
have levels proportional to their weights.
\end{cor}

For three or more cliques the pairwise ratios
$W(D_a)/W(D_b)$ must also be consistent around every cycle, which for
overlapping cliques of different weights generically fails, and then only
the level $0$ remains.

\begin{thm}[collision]
\label{thm:collision}
Let $H$ anchor distinct maximal cliques $Q_1,\ldots,Q_k$, $k\geq2$, at a
common level $\mu\neq0$, and let $E=\ker(\mu I-\hH)\cap z^\perp$. Then
\begin{enumerate}
\item[(a)] every difference $x^{Q_a}-x^{Q_b}$ lies in $E$, so $\mu$ is an
eigenvalue of $\hH$ of multiplicity at least
$d=\dim\aff\{x^{Q_1},\ldots,x^{Q_k}\}\geq1$;
\item[(b)] $\hb\perp E$: the linear form $c$ vanishes on the whole
eigenspace;
\item[(c)] $\hx^{Q_a}=\hx_p+e_a$ with a common $\hx_p\in E^\perp$ and
$e_a\in E$; if the $W(Q_a)$ are equal, no $x^{Q_a}$ lies on the secular
curve, which passes through $x^0+\hx_p$ at $\mu$.
\end{enumerate}
\end{thm}
\begin{proof}
By Lemma~\ref{lem:onsecular}, $(\mu I-\hH)\hx^{Q_a}=\hb$ for every $a$;
subtracting two of these gives (a). Then $\hb$ lies in the range of the
symmetric operator $\mu I-\hH$ on $z^\perp$, which is $E^\perp$, giving
(b). Decompose $\hx^{Q_a}$ along $E^\perp\oplus E$; the $E^\perp$
components solve the same nonsingular system on $E^\perp$ and so coincide.
If the weights are equal,
$\|\hx^{Q_a}\|^2=1/W(Q_a)-1/W(V)$ is the same for all $a$, hence so is
$\|e_a\|$; if one $e_a$ vanished all would, and all the indicators would
coincide. Since $c=0$ on $E$, the secular curve has no pole at $\mu$ and
tends to $x^0+(\mu I-\hH)^+\hb=x^0+\hx_p$ there; at any other multiplier
it cannot meet $x^{Q_a}$, whose multiplier is $\mu$ by
Lemma~\ref{lem:onsecular}.
\end{proof}

This is precisely the ``degenerate case of multiple eigenvalues with
$c_i=0$'' that \cite[Section~5]{CW} sets aside: anchoring several cliques
of equal weight forces it, whatever the remaining free entries. The
contrapositive is useful as well: if $\hb$ has a nonzero component along
$x^{Q_a}-x^{Q_b}$ for two cliques of equal weight, then $H$ does not anchor
both. Making a degenerate direction ``active'' ($c\neq0$ on it) and
anchoring the cliques that span it are mutually exclusive.

\begin{thm}[two cliques]
\label{thm:two}
Let $Q_a\neq Q_b$ be maximal cliques and let $B$ be the bipartite graph on
$D_a\cup D_b$ whose edges are the non-adjacent pairs of $G$ between $D_a$
and $D_b$. There is $H\in\Hc(G,w)$ anchoring $Q_a$ and $Q_b$ at nonzero
levels iff
\begin{enumerate}
\item[(i)] every connected component $\Gamma$ of $B$ satisfies
$W(\Gamma\cap D_a)/W(\Gamma\cap D_b)=W(D_a)/W(D_b)$, and
\item[(ii)] if some vertex outside $Q_a\cup Q_b$ is adjacent to all of
$Q_a\triangle Q_b$, then $W(Q_a)=W(Q_b)$.
\end{enumerate}
The levels then satisfy $\mu_a:\mu_b=W(D_a):W(D_b)$.
\end{thm}
\begin{proof}
Every vertex of $B$ has a neighbour in $B$: a vertex $j\in D_a$ adjacent to
all of $D_b$ would be adjacent to all of $Q_b$, contradicting maximality.
The equations (\ref{eq:system}) and the defects they involve split into
independent blocks.

{\em Outside vertices.} For $i\notin Q_a\cup Q_b$ let $A_i=D_a\setminus
N(i)$, $B_i=D_b\setminus N(i)$, $K_i=K\setminus N(i)$. The equations of $i$
are $\tau(K_i)+\tau(A_i)=\mu_a$ and $\tau(K_i)+\tau(B_i)=\mu_b$, where
$\tau(S)=\sum_{j\in S}w_j\tau_{ij}$, and their defects occur in no other
equation. By maximality $K_i\cup A_i$ and $K_i\cup B_i$ are nonempty. If
$A_i=B_i=\emptyset$ (i.e. $i$ is adjacent to all of $Q_a\triangle Q_b$),
both read $\tau(K_i)=\mu_a=\mu_b$; otherwise they are solvable for any
levels (if, say, $B_i=\emptyset\neq A_i$, set $\tau(K_i)=\mu_b$ and
$\tau(A_i)=\mu_a-\mu_b$).

{\em The two differences.} As in the proof of Theorem~\ref{thm:level} the
remaining equations are, with $\xi_{ij}=w_iw_j\tau_{ij}$ on the edges of
$B$, $\sum_j\xi_{ij}=w_i\mu_a$ for $i\in D_b$ and $\sum_i\xi_{ij}=w_j\mu_b$
for $j\in D_a$: edge values of $B$ with prescribed sums at the vertices.
The incidence matrix of a connected bipartite graph has corank one, its
left kernel spanned by the vector that is $+1$ on one side and $-1$ on the
other; so this is solvable iff every component satisfies
$\mu_aW(\Gamma\cap D_b)=\mu_bW(\Gamma\cap D_a)$. Summing over components
gives Theorem~\ref{thm:level}, and for nonzero levels the component
conditions hold iff (i) does. With (ii) for the outside vertices this
proves the claim.
\end{proof}

\begin{cor}
\label{cor:twomax}
Any two maximum weight cliques can be anchored in one wrapper at a common
positive level with all defects nonnegative.
\end{cor}
\begin{proof}
Here $W(D_a)=W(D_b)$ and the level is common. For $S\subseteq D_b$ let
$N_B(S)$ be its neighbourhood in $B$. The set $(Q_a\setminus N_B(S))\cup S$
is a clique (vertices of $S$ are adjacent to $K$, and to every vertex of
$D_a$ outside $N_B(S)$), of weight $W(Q_a)-W(N_B(S))+W(S)$, so
$W(N_B(S))\geq W(S)$ by maximality of $W(Q_a)$. By the supply--demand
theorem of Gale \cite{G57}, the transportation problem with supplies
$w_i\mu$ on $D_b$, demands $w_j\mu$ on $D_a$ and routes along the edges of
$B$ then has a nonnegative solution. For an outside vertex $i$ take
$\tau(K_i)=\mu$ and $\tau(A_i)=\tau(B_i)=0$ if $A_i$ or $B_i$ is empty
($K_i$ is then nonempty), and $\tau(K_i)=0$, $\tau(A_i)=\tau(B_i)=\mu$
otherwise, spreading each sum evenly over its non-edges.
\end{proof}

Nonnegative, however, does not mean proper:

\begin{prop}[forced blindness]
\label{prop:blind}
Let $H$ anchor $Q_a$ and $Q_b$ at a common level, and let
$i\notin Q_a\cup Q_b$ be adjacent to all of $D_b$. Then
$\sum_{j\in D_a\setminus N(i)}w_j\tau_{ij}=0$. In particular a proper
wrapper cannot anchor both if such an $i$ has a non-neighbour in $D_a$;
with nonnegative defects all these $\tau_{ij}$ vanish, so $h_{ij}=z_iz_j$ on
those non-edges: the wrapper does not distinguish them from edges.
\end{prop}
\begin{proof}
Subtract the two equations of $i$ in the proof of Theorem~\ref{thm:two}:
with $B_i=\emptyset$ they give $\tau(A_i)=\mu_a-\mu_b=0$.
\end{proof}

In dense graphs, maximum cliques that differ by a single exchange
$Q_b=Q_a-j+j'$ are common, and every outside vertex $i$ adjacent to $j'$ but
not to $j$ then makes the non-edge $\{i,j\}$ invisible, here without any sign
condition, since $D_a=\{j\}$ leaves a single defect in the sum. Once invisible
non-edges cover all the non-neighbours of some vertex in some other anchored
clique, that vertex's equation reads $0=\mu$ and the common level collapses;
Example~\ref{ex:cascade} shows this with three maximum cliques.

\section{Anchoring every maximum clique}
\label{sec:all}

\subsection{The best relaxation is the Lov\'asz number}

\begin{thm}
\label{thm:theta}
For every $H\in\Hc(G,w)$, $v(H)\leq 1/\vt$, with equality whenever
$\lmax(H)=\vt$. Hence
\[ \max_{H\in\Hc(G,w)}\ \min\{x^Tx:\ x\in\Wc^H\}=\frac{1}{\vt}. \]
\end{thm}
\begin{proof}
{\em Equality for Lov\'asz-optimal $H$.} For $x\in\Wc^H$,
$1=x^THx\leq\lmax(H)\,x^Tx=\vt\,x^Tx$, so $v(H)\geq1/\vt$; the upper bound
below then gives equality.

{\em The upper bound.} Consider the semidefinite relaxation
\[ v_{\rm SDP}(H)=\min\{\tr X:\ \langle H,X\rangle=1,\
\langle zz^T,X\rangle=1,\ X\succeq0\}. \]
If $x\in\Wc^H$ then $X=xx^T$ is feasible with $\tr X=x^Tx$, so
$v_{\rm SDP}(H)\leq v(H)$. Conversely, the feasible set is nonempty and the
trace is bounded below and coercive on it, so the minimum is attained and
the optimal set is a nonempty compact face of the feasible set. An extreme
point of it is an extreme point of the feasible set, which is cut out of
the positive semidefinite cone by two linear equations; by the rank bound
of Barvinok and Pataki \cite{Ba95,P98} it has rank $r$ with
$r(r+1)/2\leq2$, so $r\leq1$, and $r=0$ is excluded by
$\langle zz^T,X\rangle=1$. Writing it as $xx^T$ with $z^Tx=1$ (replacing $x$ by
$-x$ if necessary) gives $x\in\Wc^H$ with $x^Tx=v_{\rm SDP}(H)$. Hence
$v(H)=v_{\rm SDP}(H)$.

Now let $X^*$ be a Lov\'asz optimum. $H$ and $zz^T$ agree on the support of
$X^*$ (the diagonal and $E$), so $X=X^*/\vt$ has
$\langle H,X\rangle=\langle zz^T,X\rangle=1$ and $\tr X=1/\vt$. Thus
$v(H)=v_{\rm SDP}(H)\leq1/\vt$.
\end{proof}

So every wrapping matrix gives an upper bound $1/v(H)\geq\vt\geq\om$ on the
weighted clique number, and the Lov\'asz-optimal wrappers give the best of
them. In the sense of the relaxed least squares program they are the best
wrappers. By Lemma~\ref{lem:gauge}, $v$ is constant on gauge orbits. For
the standard wrapper of a regular graph the bound is a classical one:

\begin{cor}
\label{cor:hoffman}
Let $G$ be $d$-regular, not complete, with unit weights, and let
$\lambda_2$ be the largest eigenvalue of $A_G$ on $\one^\perp$. Then
\[ \frac{1}{v(H_0)}=\frac{n(1+\lambda_2)}{n-d+\lambda_2}, \]
Hoffman's ratio bound \cite{Ho70} on $\alpha(\bar G)=\omega(G)$ applied to
the $(n-1-d)$-regular graph $\bar G$.
\end{cor}
\begin{proof}
Here $x^0=\one/n$ and $H_0x^0=(1+d)x^0$, so $\hb=0$ and
$s=1-(1+d)/n$. As $A_G$ preserves $\one^\perp$, $\lmax(\hH_0)=1+\lambda_2$,
which is at least $1$ for a non-complete graph ($\lambda_2$ is then the
second eigenvalue of $A_G$, at least $0$ by interlacing with two
non-adjacent vertices). The program is
$\min\|\hx\|^2$ subject to $\hx^T\hH_0\hx=s$, with value
$s/(1+\lambda_2)$, so $v(H_0)=1/n+(n-1-d)/(n(1+\lambda_2))$, which
rearranges to the claim. The smallest eigenvalue of $A_{\bar G}$ is
$-1-\lambda_2$, and Hoffman's bound for $\bar G$ is
$n(1+\lambda_2)/(n-1-d+1+\lambda_2)$.
\end{proof}

With Theorem~\ref{thm:theta} this recovers Lov\'asz's inequality
$\vartheta\leq$ Hoffman's bound for regular graphs \cite[Theorem 9]{L79}:
passing from the standard to a Lov\'asz-optimal wrapper improves the relaxed
bound from Hoffman's to $\vartheta$.

\subsection{Global anchoring}

\begin{thm}
\label{thm:global}
For a clique $Q$ the following are equivalent:
\begin{enumerate}
\item[(a)] $x^Q$ is a global minimizer of (\ref{eq:v}) for some
$H\in\Hc(G,w)$;
\item[(b)] $W(Q)=\om=\vt$.
\end{enumerate}
If they hold, every Lov\'asz-optimal $H$ anchors every maximum weight clique
at level $\vt$ ($C=0$), and every maximum weight clique indicator is a
global minimizer of (\ref{eq:v}) for it. Moreover, if some $H$ anchors a
clique $Q$ at a level $\mu>0$ with $\mu\geq\lmax(\hH)$, then (b) holds.
\end{thm}
\begin{proof}
(a)$\Rightarrow$(b): $1/W(Q)=\|x^Q\|^2=v(H)\leq1/\vt$ by
Theorem~\ref{thm:theta}, so $W(Q)\geq\vt\geq\om\geq W(Q)$.
(b)$\Rightarrow$(a): for Lov\'asz-optimal $H$, $v(H)=1/\vt=1/W(Q)=\|x^Q\|^2$.
For a Lov\'asz-optimal $H$ and a maximum weight clique $Q'$, the Rayleigh
quotient of $z_{Q'}$ is $W(Q')=\vt=\lmax(H)$ (Lemma~\ref{lem:cliquevec}),
so $z_{Q'}$ is a top eigenvector, $Hz_{Q'}=\vt z_{Q'}$, which is
Lemma~\ref{lem:tau}(c) with $\mu=W(Q')=\vt$. The last claim follows from
Lemmas~\ref{lem:onsecular} and~\ref{lem:global}(a) and (a)$\Rightarrow$(b).
\end{proof}

So the relaxed program can be exact, in the sense that some wrapper makes a
clique indicator its global minimizer, exactly when $\vartheta=\omega$; and
then {\em anchoring one maximum clique globally is the same as anchoring all
of them}. On graphs with $\vartheta>\omega$, which include most hard
benchmark graphs, every anchored clique is a non-global stationary point.
This is why QUALEX-MS must examine the stationary points in all intervals of
the secular equation and not only the one beyond $\lmax$.

\begin{cor}[the hard case]
\label{cor:hard}
Let $\vt=\om$ and let $G$ have $k\geq2$ maximum weight cliques, whose
indicators span an affine space of dimension $d$. For every Lov\'asz-optimal
$H$, $\mu=\vt$ equals $\lmax(\hH)$, with multiplicity at least $d$ and with
$\hb$ orthogonal to its eigenspace. The multiplicity of $\lmax(H)$ is at
least the rank of every Lov\'asz optimum, in particular at least
$\dim\spn\{z_Q\}$ over the maximum weight cliques $Q$.
\end{cor}
\begin{proof}
All maximum cliques are anchored at $\mu=\vt$ (Theorem~\ref{thm:global}).
The restriction of $\hH$ to $z^\perp$ is a compression of $H$, so
$\lmax(\hH)\leq\lmax(H)=\vt$ by interlacing, while
Theorem~\ref{thm:collision} makes $\mu$ an eigenvalue of $\hH$ of
multiplicity at least $d$ with $\hb$ orthogonal to its eigenspace. By
(\ref{eq:compl}) the range of every Lov\'asz optimum $X$ lies in
$\ker(\vt I-H)$, and the matrices $X_Q$ of Lemma~\ref{lem:cliquevec} are
Lov\'asz optima when $W(Q)=\vt$.
\end{proof}

In trust region terms, the sphere of radius $\|\hx^Q\|$ is in the {\em hard
case} \cite{MS83}: its global maximizers form a sphere inside
$x^0+\hx_p+E$, which contains all the maximum clique indicators, none of
them on the secular curve. This is the structural reason why the
Lov\'asz-optimal wrappers, although best by Theorem~\ref{thm:theta}, are
useless for recognizing a clique through the secular equation when there are
several maximum cliques. When $\vartheta>\omega$ the Lov\'asz-optimal wrappers
are degenerate for a different reason, the generic multiplicity of optimal
eigenvalues \cite{O92,P98}; Section~\ref{sec:comp} shows both.

\subsection{Colourings, holes and perfect graphs}

In this subsection all weights are $1$, so $z=\one$ and $\hH$, $\tau$ are
those of an unweighted graph.

\begin{prop}[the colouring wrapper]
\label{prop:colour}
Let $\chi(G)=\omega(G)=\omega$, with colour classes $S_1,\ldots,S_\omega$.
Then
\[ H_\chi=\omega I-\Bigl(\omega\sum_{c=1}^{\omega}\one_{S_c}\one_{S_c}^T
-\one\one^T\Bigr) \]
lies in $\Hc(G,\one)$, with defects $\tau_{ij}=\omega$ inside colour
classes and $\tau_{ij}=0$ between them, has $\lmax(H_\chi)=\omega$ and
anchors every maximum clique at level $\omega$ ($C=0$).
\end{prop}
\begin{proof}
The diagonal of $H_\chi$ is $\omega-(\omega-1)=1$; between classes the
entries are $1$, which covers all edges; inside a class (non-adjacent
pairs) they are $1-\omega$. The bracket is positive semidefinite by the
Cauchy--Schwarz inequality,
$\omega\sum_c(\one_{S_c}^Tx)^2\geq(\sum_c\one_{S_c}^Tx)^2$, so
$H_\chi\preceq\omega I$, and $\lmax(H_\chi)=\omega$ by
Lemma~\ref{lem:cliquevec}. A maximum clique $Q$ meets every class in one
vertex; for $i\notin Q$ in class $c$, $\sum_{j\in Q}h_{ij}=(1-\omega)+
(\omega-1)=0$, so $C=0$.
\end{proof}

The colouring wrapper sees only the complete $\omega$-partite graph on the
colour classes: the indicator of every transversal of the colouring lies on
$\Wc^{H_\chi}$ and has norm $1/\sqrt{\omega}$, exactly like a maximum
clique, although most transversals are not cliques. Such pseudo-cliques are
what the strict inequality of \cite[Proposition~2]{CW} rules out; the
colouring wrapper is improper as soon as two colour classes are not
completely joined.

\begin{thm}
\label{thm:holes}
Let $G$ be an odd hole $C_{2k+1}$ or an odd antihole $\bar C_{2k+1}$,
$k\geq2$, with unit weights. Every $H\in\Hc(G,\one)$ that anchors all
maximum cliques of $G$ does so at level $0$, and then $H=\one\one^T$.
\end{thm}
\begin{proof}
All maximum cliques share one level $\mu$ (Corollary~\ref{cor:equal}).
Vertices are taken modulo $n=2k+1$.

{\em Odd holes.} The maximum cliques are the edges $\{a,a+1\}$. For
$Q=\{a,a+1\}$ the outside vertex $a+2$ misses only $a$ and $a-1$ misses
only $a+1$, so every non-edge at distance $2$ has $\tau=\mu$; a vertex $v$
with $a+3\leq v\leq a-2$ misses both, giving
$\tau_{a,v}+\tau_{a+1,v}=\mu$. With $v=a+j+1$ this reads
$\tau_{a,a+j+1}=\mu-\tau_{a+1,a+j+1}$, so by induction on $j$ every pair
$\{b,b+j\}$ has $\tau=\mu$ for even $j$ and $\tau=0$ for odd $j$,
$2\leq j\leq n-2$. The pair $\{b,b+2\}$ is also the pair
$\{b+2,(b+2)+(n-2)\}$ with $n-2$ odd, so $\mu=0$, and then every defect is
$0$.

{\em Odd antiholes.} The non-edges are the pairs $\{b,b+1\}$; write
$\tau_b$ for their defects. The maximum cliques of $\bar C_{2k+1}$ are the
stable sets of size $k$ of $C_{2k+1}$, namely
$Q_a=\{a,a+2,\ldots,a+2k-2\}$ (consecutive elements differ by $2$ except for
one gap of $3$). The non-neighbours of an outside vertex $v$ in $Q_a$ are
its neighbours on the cycle that belong to $Q_a$. For $v=a+2k-1$ this is
$\{a+2k-2\}$ and for $v=a+2k$ it is $\{a\}$, so $\tau_{a+2k-2}=\mu$ and
$\tau_{a-1}=\mu$ for all $a$: every defect equals $\mu$. For $v=a+1$ (which
is outside $Q_a$ since $k\geq2$) it is $\{a,a+2\}$, giving
$\tau_a+\tau_{a+1}=\mu$, i.e. $2\mu=\mu$. So $\mu=0$ and all defects vanish.
\end{proof}

Even holes behave differently. The same induction shows that $H$ anchors all
maximum cliques of $C_{2k}$ ($k\geq2$) iff $\tau=\mu$ at even distances and
$\tau=0$ at odd distances, which is consistent because $j$ and $2k-j$ have
the same parity. At $\mu=2$ this is the colouring wrapper of the
bipartition, so it is Lov\'asz-optimal ($\vartheta=\omega=2$). For $k\geq3$
the odd-distance non-edges are invisible to {\em every} wrapper anchoring
all maximum cliques, hence to every Lov\'asz-optimal wrapper. For $C_6$ the
three long diagonals $\{i,i+3\}$ then satisfy $x^THx=1$ for
$x=(e_i+e_{i+3})/2\in\Delta_n$, which has $x^Tx=1/2=1/\omega$: they are
global minimizers of the program of \cite[Theorem~1]{CW} although they are not
cliques. The Lov\'asz-optimal wrappers of $C_6$ are necessarily improper.

\begin{thm}
\label{thm:perfect}
For a graph $G$ with unit weights the following are equivalent:
\begin{enumerate}
\item[(a)] $G$ is perfect;
\item[(b)] for every induced subgraph $G'$ of $G$ some $H\in\Hc(G',\one)$
anchors all maximum cliques of $G'$ at a common nonzero level.
\end{enumerate}
When (a) holds, $H$ in (b) can be taken Lov\'asz-optimal for $G'$.
\end{thm}
\begin{proof}
(a)$\Rightarrow$(b): induced subgraphs of perfect graphs are perfect, so
$\chi(G')=\omega(G')$, and Proposition~\ref{prop:colour} applies.
(b)$\Rightarrow$(a): if $G$ is not perfect, by the strong perfect graph
theorem \cite{CRST06} it has an induced odd hole or odd antihole, for which
(b) fails by Theorem~\ref{thm:holes}.
\end{proof}

Condition (b) is linear (no semidefiniteness is involved), but it must be
required hereditarily: an imperfect graph with a unique maximum clique
satisfies it for $G'=G$, and so do the graphs of Section~\ref{sec:comp} with
$\vartheta>\omega$ whose maximum cliques can all be anchored. One direction
holds without heredity:

\begin{cor}
\label{cor:certificate}
If the maximum weight cliques of $G$ cannot all be anchored at a common
nonzero level, then $\vt>\om$.
\end{cor}
\begin{proof}
If $\vt=\om$, a Lov\'asz-optimal wrapper anchors them all at level $\vt$
(Theorem~\ref{thm:global}).
\end{proof}

\begin{example}
\label{ex:cascade}
The unweighted random graph \gr{g200_80_0} ($n=200$, $p=0.8$, $\omega=26$)
has three maximum cliques. $Q_2=Q_1-14+53$ (vertex labels of the DIMACS
file), and $Q_3$ shares only four vertices with them. The vertices 95, 113,
158 and 176 of $Q_3$ are adjacent to 53 but not to 14, so by
Proposition~\ref{prop:blind} (with $D_a=\{14\}$, a single vertex, so
without any sign condition) $\tau_{i,14}=0$ for each of them. But these four
are exactly the non-neighbours of vertex 14 in $Q_3$, so the equation of 14
for $Q_3$ reads $0=\mu$. The system has 418 distinct equations in 1693
defects and is still inconsistent; this certificate combines nine of them.
\end{example}

\subsection{Quasigroup completion: why the hard case does not help}
\label{sec:qcp}

Busygin and Pasechnik \cite{BP03} reduce the quasigroup completion problem
(QCP) to the maximum independent set problem. For a partial latin square of
order $n$ with $h$ empty cells (holes), the {\em QCP graph} $G$ has as
vertices the triples $(i,j,k)$ such that $(i,j)$ is a hole and the symbol $k$
occurs neither in row $i$ nor in column $j$, and two triples are adjacent when
they agree in two coordinates; it is the subgraph of the Hamming graph
$H(3,n)$ induced by the non-neighbours of the prefilled triples. Its
independent sets have at most $h$ vertices, and the independent sets of size
$h$ are exactly the completions of the square \cite[Lemma~2]{BP03}. The
candidates of one hole form a clique, so $\overline{\chi}(G)\leq h$
\cite[Lemma~3]{BP03}, and deciding whether $\alpha(G)=\overline{\chi}(G)$ is
NP-hard even when this clique partition is given \cite[Theorem~1]{BP03}.
Hence every completable instance has
\[ \alpha(G)=\vartheta(G)=\overline{\chi}(G)=h, \]
with $\vartheta(G)$ the bound (\ref{eq:thetaprimal}) for the graph $\bar G$,
and the hard instances of QCP are among these or look like them to
$\vartheta$. The independent sets of $G$ are the cliques of $\bar G$, so we
use the wrapping matrices of $\bar G$ with unit weights; defects now sit on
the edges of $G$, and the standard wrapping matrix of $\bar G$ is
$J-A_G$.

Besides the cells, $G$ has two more partitions into cliques: the
row--symbol lines $\{(i,\cdot,k)\}$ and the column--symbol lines
$\{(\cdot,j,k)\}$. Two distinct triples share at most one line (sharing two
coordinates of two different kinds would make them equal).

\begin{prop}
\label{prop:qcp}
Let $G$ be the QCP graph of a completable instance with $h$ holes, and let
$\mathcal{P}_1,\mathcal{P}_2,\mathcal{P}_3$ be its partitions into cells,
row--symbol lines and column--symbol lines, each of $h$ cliques. Then:
\begin{enumerate}
\item[(a)] $H=J-\frac h3A_G$, the standard wrapping matrix of $\bar G$ up to
gauge, has $\lmax(H)=h=\vartheta(G)$; it is Lov\'asz-optimal and anchors
every completion $S$ at level $h$ ($C=0$);
\item[(b)] its eigenspace for $h$ is
$L=\{x:\ \one_C^Tx \mbox{ is the same for all } C\in\mathcal{P}_t,\
t=1,2,3\}$;
\item[(c)] the global minimizers of (\ref{eq:v}) for $H$ are the $x\in L$
with $\one^Tx=1$ and $x^Tx=1/h$, and the nonnegative ones among them are
exactly the indicators $\one_S/h$ of the completions $S$.
\end{enumerate}
\end{prop}
\begin{proof}
For $t=1,2,3$ let $K_t=h\sum_{C\in\mathcal{P}_t}\one_C\one_C^T-J$, which is
positive semidefinite by the Cauchy--Schwarz inequality as in
Proposition~\ref{prop:colour}, and $H_t=hI-K_t$, the colouring wrapper of
$\bar G$ for the colouring $\mathcal{P}_t$. $H_t$ has $1-h$ on the pairs that
share a line of $\mathcal{P}_t$ and $1$ elsewhere. Every edge of $G$ shares
exactly one line, so $(H_1+H_2+H_3)/3$ has $1$ on the diagonal and on the
non-adjacent pairs of $G$ and $1-h/3$ on its edges: it is $H$, with defect
$h/3$ on every non-edge of $\bar G$, i.e. $(J-A_G)_t$ with $1+t=h/3$ in the
notation of Lemma~\ref{lem:gauge}. Moreover $H=hI-\frac13(K_1+K_2+K_3)
\preceq hI$, and $\ker(hI-H)=\bigcap_t\ker K_t$, where $\ker K_t$ consists
of the vectors whose sums over the $h$ classes of $\mathcal{P}_t$ are equal
(the equality case of the Cauchy--Schwarz inequality). This is (b). A
completion meets every line exactly once, so $\one_S\in L$ and
$H\one_S=h\one_S$, which gives $\lmax(H)=h$, the anchoring at level $h$
(Lemma~\ref{lem:tau}(c)), and $\vartheta(G)\leq h=\alpha(G)\leq
\vartheta(G)$.

For (c), $v(H)=1/h$ by Theorem~\ref{thm:theta}. A point of $\Wc^H$ has
$1=x^THx\leq h\,x^Tx$, with equality iff $x\in L$; so the global minimizers
are the points of $L$ with $\one^Tx=1$ and $x^Tx=1/h$ (on $L$,
$x^THx=h\,x^Tx$). Let such an $x$ be nonnegative. Each partition has $h$
classes covering $V$, so all its class sums equal $1/h$, and
\[ x^Tx=\sum_{\rm holes}\ \sum_k x_{ijk}^2\ \leq\ \sum_{\rm holes}
\Bigl(\sum_k x_{ijk}\Bigr)^2=\frac1h, \]
with equality iff every hole carries a single candidate, with value $1/h$.
Those candidates meet every row--symbol and column--symbol line once, so they
form a completion.
\end{proof}

So the standard wrapper that QUALEX-MS builds for the independent sets of a
completable QCP graph is Lov\'asz-optimal, its trust region program at the
radius of the completions is in the hard case with the degenerate
eigenspace $L\cap\one^\perp$, and the completions are exactly the points of
the sphere of global minimizers that lie in the nonnegative orthant, i.e.
the points of largest norm of the polytope $\{x\in L:\ x\geq0,\ \one^Tx=1\}$
of fractional completions. The hard case therefore carries all the
information about the completions, but finding a nonnegative point on that
sphere is QCP itself, which is NP-complete \cite{Co84}. The colouring wrapper
of the cells alone (Proposition~\ref{prop:colour} for $\mathcal{P}_1$) is
Lov\'asz-optimal as well but blind to every row and column conflict: by the
same argument every transversal, one candidate for each hole, is a global
minimizer, and its top eigenspace has dimension $|V|-h+1$. For the empty
square, $G=H(3,n)$, $L$ is the sum $V_0\oplus V_3$ of the eigenspaces of the
Hamming scheme, of dimension $(n-1)^3+1$, and (a) is the tightness of
Hoffman's bound $\alpha(H(3,n))=n^2$ (Corollary~\ref{cor:hoffman}).

Every Lov\'asz-optimal wrapper has at least the rank of every Lov\'asz
optimum as the multiplicity of its top eigenvalue (Corollary~\ref{cor:hard}),
and that rank is at least the dimension spanned by the completions.
Section~\ref{sec:qwh} measures where between that span and $\dim L$ the
Lov\'asz optimum falls on random instances.

\subsection{SAT01: exact covers with contradictions}
\label{sec:sat01}

Busygin's SAT01 \cite{Bu07} puts QCP into a general weighted setting. An
instance SAT01$(A,\Gamma)$ is a $0$-$1$ matrix $A$ of size $m\times n$
without zero columns and a graph $\Gamma$ on the $n$ variables, and it asks
for $x\in\{0,1\}^n$ with $Ax=\one$ whose support is independent in $\Gamma$.
We identify each equation with the set $e\subseteq V$ of its variables and
write $\mathcal E$ for the $m$ equations. $\Gamma$ contains every pair of
variables that share an equation and possibly further pairs; all its edges
are called {\em contradictions}. With the weights $w=A^T\one$, the number of
equations of each variable, an independent set of $\Gamma$ meets every
equation at most once, so its weight is at most $m$, with equality iff it is
the support of a solution \cite[Lemma~1]{Bu07}. The equations are $m$
cliques of $\Gamma$ that cover every vertex $j$ exactly $w_j$ times, so
$\vartheta(\Gamma,w)\leq m$ \cite[Theorem~1]{Bu07}. Many NP problems reduce to
SAT01 in a natural way (CNF satisfiability, Hamilton cycles, graph and
subgraph isomorphism, QCP, the 15-puzzle, the Hanoi tower, blocksworld
planning, factoring). A satisfiable instance
has $\alpha=\vartheta=m$, so $\vartheta<m$ refutes an instance in polynomial
time, and only the unsatisfiable instances with $\alpha<\vartheta=m$ escape
it.

In the notation of this note $G=\bar\Gamma$: two variables are adjacent iff
they do not contradict, the solutions are the cliques of weight $m$ (maximal
cliques, \cite[Claim~1]{Bu07}), $\om=\alpha(\Gamma,w)$,
$\vt=\vartheta(\Gamma,w)$, and the non-edges of $G$ are the contradictions.
The QCP graph of Section~\ref{sec:qcp} is the instance with the cells,
row--symbol lines and column--symbol lines as equations, all weights $3$,
$m=3h$, and no contradictions besides the equations. Let
$D=\operatorname{diag}(w)$, so that $z=D^{1/2}\one$, and $M=A^TA$: $M_{ij}$
counts the equations shared by $i$ and $j$, and $M_{jj}=w_j$.

\begin{prop}
\label{prop:sat01}
Let
\[ H_A = zz^T + m\,\bigl(I - D^{-1/2}MD^{-1/2}\bigr). \]
\begin{enumerate}
\item[(a)] $H_A\in\Hc(G,w)$, with defect $\tau_{ij}=mM_{ij}/(w_iw_j)$ on
every contradiction; in particular $\tau_{ij}=0$ on the contradictions that
share no equation.
\item[(b)] $mI-H_A=D^{-1/2}A^T(mI-\one\one^T)AD^{-1/2}\succeq0$, so
$\lmax(H_A)\leq m$, and $\ker(mI-H_A)=L_A:=D^{1/2}\{u:\ Au\in\R\one\}$.
\item[(c)] If $\vt=m$, $H_A$ is Lov\'asz-optimal; conversely, if $L_A\neq0$
(e.g. when $n>m$) and $H_A$ is Lov\'asz-optimal, $\vt=m$. If the instance is
satisfiable, $\vt=m$ and $H_A$ anchors every solution $S$ at level $m$
($C=0$).
\item[(d)] If $\vt=m$, the global minimizers of (\ref{eq:v}) for $H_A$ are
the $x\in L_A$ with $z^Tx=1$ and $x^Tx=1/m$, and the nonnegative ones among
them are exactly the indicators $x^S=z_S/m$ of the exact covers $S$ of $A$
(the $0$-$1$ solutions of $Ax=\one$), whether or not they respect the
contradictions outside the equations.
\end{enumerate}
\end{prop}
\begin{proof}
(a) As $A$ is $0$-$1$, $M_{jj}=w_j$ and the diagonal of $H_A$ is
$w_j+m(1-1)=w_j$. Off the diagonal
$(H_A)_{ij}=z_iz_j-mM_{ij}/(z_iz_j)=z_iz_j\bigl(1-mM_{ij}/(w_iw_j)\bigr)$.
Adjacent vertices of $G$ share no equation, so there $(H_A)_{ij}=z_iz_j$.

(b) As $D^{-1/2}A^T\one=D^{-1/2}w=z$,
\[ mI-H_A=m\,D^{-1/2}MD^{-1/2}-zz^T
   =D^{-1/2}A^T\bigl(mI-\one\one^T\bigr)AD^{-1/2}. \]
Since $mI-\one\one^T$ is positive semidefinite with kernel $\R\one$, $y$ lies
in $\ker(mI-H_A)$ iff $AD^{-1/2}y\in\R\one$.

(c) By (\ref{eq:thetaprimal}) and (b), $\vt\leq\lmax(H_A)\leq m$, and
$\lmax(H_A)=m$ when $L_A\neq0$; both implications follow. A solution $S$ is
a clique of weight $m$, so
$\vt\geq m$ by Lemma~\ref{lem:cliquevec}, and $AD^{-1/2}z_S=A\one_S=\one$,
so $z_S\in L_A$ and $H_Az_S=mz_S$, which is the anchoring at level
$m=W(S)$ by Lemma~\ref{lem:tau}(c).

(d) $\lmax(H_A)=m=\vt$, so $v(H_A)=1/m$ by Theorem~\ref{thm:theta}, and a
point of $\Wc^{H_A}$ has $1=x^TH_Ax\leq m\,x^Tx$, with equality iff
$x\in L_A$; this gives the global minimizers. Write such an $x$ as
$D^{1/2}u$ with $Au=t\one$. Then $z^Tx=w^Tu=\one^TAu=tm$, so $t=1/m$, and
$p=mu$ satisfies $Ap=\one$. If $x\geq0$ then $p\geq0$, hence $p\leq\one$,
since every variable lies in an equation on which $p$ sums to $1$, and
\[ m^2\,x^Tx=\sum_j w_jp_j^2\ \leq\ \sum_j w_jp_j=\one^TAp=m, \]
with equality iff $p\in\{0,1\}^n$. So $x^Tx=1/m$ iff $p=\one_S$ for an exact
cover $S$, and then $x=D^{1/2}\one_S/m=x^S$.
\end{proof}

For QCP every pair shares at most one line, so $M=3I+A_\Gamma$ and
$H_A=3J+3hI-h(3I+A_\Gamma)=3\bigl(J-\frac h3A_\Gamma\bigr)$, three times the
wrapper of Proposition~\ref{prop:qcp} (whose weights are $1$), with
$L_A=L$; that proposition is the case of (a)--(d) in which every exact cover
is a solution. In general $H_A$ sees only the equations: it is blind to
every contradiction that shares no equation, and the nonnegative points of
its sphere of global minimizers include the exact covers that violate such
contradictions. It is always one of the Lov\'asz-optimal wrappers when
$\vartheta=m$, but on an unsatisfiable instance with $\vartheta<m$ it is
not. Which instances these are is decided by the following relaxation of the
solutions.

\begin{thm}
\label{thm:vector}
$\vt=m$ iff SAT01$(A,\Gamma)$ has a {\em vector solution}: vectors
$q_1,\ldots,q_n$ in some $\R^r$ and a unit vector $d$ with
\[ q_i^Tq_j=0 \mbox{ for every contradiction } \{i,j\}, \qquad
   \sum_{j\in e}q_j=d \mbox{ for every equation } e\in\mathcal E. \]
The Lov\'asz optima are then exactly the matrices
$X=\bigl(z_iz_j\,q_i^Tq_j/m\bigr)$ of the vector solutions, so
$\rank X=\dim\spn\{q_j\}$, and the vector solutions of dimension one are the
solutions, $q_j=x_jd$. Every vector solution gives a fractional exact cover
$p_j=\|q_j\|^2=d^Tq_j$, with $Ap=\one$, $p\geq0$ and $\sum_{j\in K}p_j\leq1$
for every clique $K$ of $\Gamma$.
\end{thm}
\begin{proof}
Variables that share an equation contradict, so the vectors of one equation
are pairwise orthogonal and $\sum_{j\in e}\|q_j\|^2=\|d\|^2=1$.

Given a vector solution, let $v_j=z_jq_j/\sqrt m$ and let $X$ be their Gram
matrix. It is positive semidefinite and vanishes on the contradictions, i.e.
off $E$, and
$\tr X=\frac1m\sum_jw_j\|q_j\|^2=\frac1m\sum_{e\in\mathcal E}\sum_{j\in e}
\|q_j\|^2=1$. Moreover
$\sum_jz_jv_j=\frac1{\sqrt m}\sum_jw_jq_j=\frac1{\sqrt
m}\sum_{e\in\mathcal E}\sum_{j\in e}q_j=\sqrt m\,d$, so
$\langle zz^T,X\rangle=\|\sum_jz_jv_j\|^2=m$. Hence $\vt\geq m$, and $\vt=m$
by Proposition~\ref{prop:sat01}(b); $X$ is a Lov\'asz optimum.

Conversely, let $\vt=m$ and let $X$ be a Lov\'asz optimum. $H_A$ is
Lov\'asz-optimal (Proposition~\ref{prop:sat01}(c)), so by (\ref{eq:compl})
the range of $X$ lies in $L_A$. Factor $X=VV^T$ with $V$ of full column rank
and let $v_j$ be the rows of $V$. The columns of $V$ lie in the range of $X$,
so every column of $D^{-1/2}V$ is a $u$ with $Au\in\R\one$, and the vectors
$v_j/z_j$ have one and the same sum $t$ over every equation. Then
$\sum_jz_jv_j=\sum_jw_j\,v_j/z_j=\sum_{e\in\mathcal E}\sum_{j\in e}v_j/z_j
=mt$, and $m=\langle zz^T,X\rangle=\|mt\|^2$ gives $\|t\|=1/\sqrt m$. So
$q_j=\sqrt m\,v_j/z_j$ and $d=\sqrt m\,t$ form a vector solution:
$q_i^Tq_j=mX_{ij}/(z_iz_j)=0$ on the contradictions. The two constructions
are inverse to each other up to an orthogonal transformation of $\R^r$.

If the $q_j$ span a line, $q_j=\lambda_jd$, with $\lambda_i\lambda_j=0$ on the
contradictions and $\sum_{j\in e}\lambda_j=1$; at most one $\lambda_j$ per
equation is nonzero, so $\lambda\in\{0,1\}^n$ is a solution. Finally, for
$j$ in an equation $e$, $d^Tq_j=\sum_{i\in e}q_i^Tq_j=\|q_j\|^2$, so
$p_j=\|q_j\|^2=d^Tq_j\geq0$ does not depend on $e$ and sums to $\|d\|^2=1$
over every equation. On a clique $K$ of $\Gamma$ the nonzero $q_j$ are
pairwise orthogonal, and Bessel's inequality gives
$\sum_{j\in K}p_j=\sum_{j\in K,\,q_j\neq0}(d^Tq_j/\|q_j\|)^2\leq\|d\|^2=1$.
\end{proof}

A vector solution is the semidefinite relaxation of the $0$-$1$ system: each
$x_j$ becomes a vector $q_j$, $x_ix_j=0$ becomes orthogonality and
$\sum_{j\in e}x_j=1$ becomes $\sum_{j\in e}q_j=d$. So the gap $\alpha<\vartheta=m$
of \cite{Bu07} is the existence of vector solutions without solutions. By the
last statement it implies that the fractional exact cover polytope
$\{p\geq0:\ Ap=\one,\ p(K)\leq1$ for all cliques $K$ of $\Gamma\}$ is nonempty,
i.e. $\kappa(\Gamma,w)=m$ in the notation of \cite{Bu07}.

SAT01 propagation \cite[Section~3]{Bu07} reduces instances by unit
propagation and clique analysis and adds edges to $\Gamma$ by contradiction
analysis: $j$ and $k$ contradict if some equation consists of variables that
contradict $j$ or $k$. Every solution survives, so $\alpha$ is unchanged, but
a vector solution need not have $q_j^Tq_k=0$ for such a derived pair, so the
new edges can destroy all vector solutions and push $\vartheta$ below $m$.
Conversely, propagation can stop short of refuting an instance with
$\vartheta<m$. Section~\ref{sec:sat01comp} shows that on
factoring instances the two refute different primes, and that their
combination refutes primes that neither refutes alone.

\subsection{Where the hardness lies}
\label{sec:hardness}

Propositions~\ref{prop:qcp} and~\ref{prop:sat01} raise a question. A
Lov\'asz-optimal wrapper is computable in polynomial time, to any precision,
by semidefinite programming. If it is proper, \cite[Propositions~1
and~2]{CW} make the nonnegative points of its sphere of global minimizers
clique indicators, and a SAT01 instance with $\vartheta=m$ is then
satisfiable iff that sphere has a nonnegative point. So either deciding
whether such a sphere has a nonnegative point is hard, or the
Lov\'asz-optimal wrappers of SAT01 instances are in general not proper. Both
hold, but only the first is essential: it holds already where the
Lov\'asz-optimal wrapper is proper and needs no semidefinite program.

\begin{prop}
\label{prop:proper}
Let $H\in\Hc(G,w)$ be Lov\'asz-optimal, $T=\ker(\vt I-H)$, and
\[ \Sigma_H=\{x\in T:\ z^Tx=1,\ x^Tx=1/\vt\}, \qquad
   P_H=\{x\in T:\ z^Tx=1,\ x\geq0\}. \]
\begin{enumerate}
\item[(a)] $\Sigma_H$ is the set of global minimizers of (\ref{eq:v}), a
sphere (possibly a single point) in the affine space
$\{x\in T:\ z^Tx=1\}$, centred at its point of least norm.
\item[(b)] If the defects of $H$ are nonnegative, $x^Tx\leq1/\vt$ on the
polytope $P_H$, so the nonnegative points of $\Sigma_H$ are the points of
$P_H$ of largest norm, and they are vertices of $P_H$.
\item[(c)] If $H$ is proper, the nonnegative points of $\Sigma_H$ are exactly
the indicators $x^Q$ of the cliques with $W(Q)=\vt$; in particular
$\Sigma_H$ has a nonnegative point iff $\om=\vt$.
\end{enumerate}
\end{prop}
\begin{proof}
(a) By Theorem~\ref{thm:theta}, $v(H)=1/\vt$, and a point of $\Wc^H$ has
$1=x^THx\leq\vt\,x^Tx$, with equality iff $x\in T$. Conversely a point of
$T$ with $z^Tx=1$ and $x^Tx=1/\vt$ has $x^THx=\vt\,x^Tx=1$. Within the
affine space $x^Tx$ is the squared distance from the origin, so the points
of norm $1/\sqrt{\vt}$ form a sphere about its point of least norm.

(b) $P_H$ lies in the bounded set $\{x\geq0:\ z^Tx=1\}$. For $x$ there,
$x^THx\leq1$ by \cite[Proposition~1]{CW}, and on $T$, $x^THx=\vt\,x^Tx$.
A strictly convex function attains its maximum over a polytope only at
vertices.

(c) Let $x\in\Sigma_H$ be nonnegative. Then $x^THx=1$, so its support is a
clique $Q$ by \cite[Proposition~2]{CW}, and by the Cauchy--Schwarz
inequality
\[ 1=z^Tx=\sum_{i\in Q}z_ix_i\leq W(Q)^{1/2}(x^Tx)^{1/2}
   =\bigl(W(Q)/\vt\bigr)^{1/2}\leq\bigl(\om/\vt\bigr)^{1/2}\leq1. \]
Equality throughout gives $W(Q)=\om=\vt$ and $x$ proportional to $z_Q$,
i.e. $x=x^Q$. Conversely, if $W(Q)=\vt$ then $x^Q$ is a global minimizer
(Theorem~\ref{thm:global}), and it is nonnegative.
\end{proof}

So for a proper Lov\'asz-optimal wrapper, $\om=\vt$ is decided by
maximizing the norm over the polytope $P_H$. Everything convex is computed
in polynomial time: $\vt$, $H$ and $T$ by the semidefinite program, $P_H$ by
linear algebra. What remains is to find a vertex of $P_H$ of largest norm,
the maximization of a convex function over a polytope, which is NP-hard in
general. In terms of Corollary~\ref{cor:hard} and Theorem~\ref{thm:vector}
it is the step from a Lov\'asz optimum to a rank-one Lov\'asz optimum. That
this step carries the hardness, and not a lack of properness, is shown by
the instances whose contradictions all lie in equations:

\begin{cor}
\label{cor:exactcover}
For every SAT01 instance, the nonnegative points of the sphere
$\Sigma_A=\{x\in L_A:\ z^Tx=1,\ x^Tx=1/m\}$, which linear algebra computes
from $A$, are exactly the indicators $x^S$ of the exact covers $S$ of $A$.
Deciding whether $\Sigma_A$ has a nonnegative point is therefore the exact
cover problem, which is NP-complete \cite{Ka72}. If every contradiction
lies in an equation, $H_A$ is proper, and it is Lov\'asz-optimal with
$\Sigma_{H_A}=\Sigma_A$ whenever $\vartheta=m$, in particular on every
satisfiable instance.
\end{cor}
\begin{proof}
The last part of the proof of Proposition~\ref{prop:sat01}(d) uses only
$x\in L_A$, $z^Tx=1$, $x^Tx=1/m$ and $x\geq0$. By
Proposition~\ref{prop:sat01}(a) the defect $mM_{ij}/(w_iw_j)$ of a
contradiction is positive iff it shares an equation, and by (b) and (c)
$H_A$ is Lov\'asz-optimal with $\ker(mI-H_A)=L_A$ when $\vartheta=m$.
\end{proof}

For QCP this is Proposition~\ref{prop:qcp} and the remark after it; the
corollary carries it over to every exact cover problem. (Whether $\vartheta$
equals $m$ exactly is a question of precision that the corollary does not
need: on the satisfiable instances it does.) With contradictions outside the
equations properness fails as well, and the optimal face can show it:

\begin{prop}
\label{prop:improper}
Let $X$ be a Lov\'asz optimum with range $R$. If the sphere
$\{x\in R:\ z^Tx=1,\ x^Tx=1/\vt\}$ has a nonnegative point whose support is
not a clique of $G$, no Lov\'asz-optimal wrapping matrix is proper, and
every Lov\'asz-optimal one with nonnegative defects has defect $0$ on every
non-adjacent pair inside that support.
\end{prop}
\begin{proof}
By (\ref{eq:compl}), $R\subseteq\ker(\vt I-H)$ for every Lov\'asz-optimal
$H$, so such a point $x$ has $x^THx=\vt\,x^Tx=1$. With $z^Tx=1$,
\[ x^THx=(z^Tx)^2-\sum\tau_{ij}z_iz_jx_ix_j
        =1-\sum\tau_{ij}z_iz_jx_ix_j, \]
the sums running over the ordered non-adjacent pairs $i\neq j$, so the last
sum vanishes. If the defects are nonnegative, so are its terms, and they all
vanish; a proper $H$ would then make the support of $x$ a clique.
\end{proof}

The Lov\'asz optima of Section~\ref{sec:sat01comp} supply such points,
numerically. On propagated Hamilton cycle instances, alternating projections
between the sphere of Proposition~\ref{prop:improper} and the nonnegative
orthant reached points whose negative part was below $10^{-4}$ of their
norm and that are not solutions: 18 such points on the cubic graph with 20
vertices, of which at most 2 are solutions, and 4 on the dodecahedron, none
of them a solution (\texttt{bench/jam.py}, Section~4 of
\texttt{bench/runs/jam/jam-2026-10-03.txt}). So these instances have no
proper Lov\'asz-optimal wrapper. Properness fails on the contradictions
outside the equations, to which $H_A$ is blind
(Proposition~\ref{prop:sat01}(a)), but by Corollary~\ref{cor:exactcover}
this is not what keeps SAT01 hard.

Searching for nonnegative points of the sphere is thus a heuristic for an
NP-hard problem. QUALEX-MS's Douglas--Rachford stage (\texttt{QMS\_DR}) drives
the corners of each degenerate sphere of stationary points towards the
nonnegative orthant, and on $H_A$ it can use instead a greedy projection
that also respects the contradictions outside the equations
(\texttt{sat01qms -D}). On the propagated satisfiable instances where the
greedy stage of QUALEX-MS fails, it found solutions of the factoring
instances $14471$, $32399$ and $60491$, whose degenerate top clusters have
dimension 148 to 235, and of the cubic graph with 30 vertices (dimension
580), none of which QUALEX-MS solves without it. It reached 911 of $m=912$
on $1000001=101\cdot9901$ (dimension 471) and found nothing on the cubic
graphs with 40 and 50 vertices (dimensions 1178 and 1828)
(\texttt{bench/sat01.py --qms}, \texttt{bench/runs/sat01/qms.tsv}).

\subsection{A second wrapper surface}
\label{sec:surface}

A solution of a SAT01 instance lies not only on $\Sigma_A$ but, being a
clique indicator, on the surface $\Wc^H$ of every $H\in\Hc(G,w)$
(Section~\ref{sec:setting}). A proper $H$ sees the contradictions that $H_A$
does not:

\begin{prop}
\label{prop:threesets}
For every SAT01 instance and every proper $H\in\Hc(G,w)$, the nonnegative
points of $\Sigma_A\cap\Wc^H$ are exactly the indicators $x^S$ of the
solutions $S$.
\end{prop}
\begin{proof}
By Corollary~\ref{cor:exactcover} the nonnegative points of $\Sigma_A$ are
the $x^S$ of the exact covers $S$. For $x\geq0$ with $z^Tx=1$,
$x^THx=1-\sum\tau_{ij}z_iz_jx_ix_j$ over the ordered non-adjacent pairs, as
in the proof of Proposition~\ref{prop:improper}, and with all defects
positive this is $1$ iff the support of $x$ contains no non-adjacent pair.
So $x^S\in\Wc^H$ iff $S$ contains no contradiction, i.e. iff $S$ is a
solution.
\end{proof}

With the standard wrapper $H_0$, therefore, the sphere, the orthant and
$\Wc^{H_0}$ meet exactly in the solutions, while the sphere and the orthant
alone meet in the exact covers. (When every contradiction lies in an
equation, every exact cover is a solution and the surface removes nothing.)
Projecting onto a wrapper surface is the relaxed program with its origin
moved:

\begin{lem}
\label{lem:surfproj}
Let $H\in\Hc(G,w)$ and $y\in\R^n$, let $\hH=U{\rm diag}(\lambda_i)U^T$ and
$c=U^T\hb$ as in Section~\ref{sec:setting}, $\beta=U^TPy$, and let $J$ be
the open interval of the $t$ with $1+t\lambda_i>0$ for every $i$, on which
$I+t\hH$ is positive definite on $z^\perp$ (between $-1/\lmax(\hH)$ and
$-1/\lmin(\hH)$ when $\lmax(\hH)>0>\lmin(\hH)$). For $t\in J$ let
\[ u_i(t)=\frac{\beta_i-tc_i}{1+t\lambda_i}, \qquad
   \phi(t)=\sum_i(\lambda_iu_i+2c_i)\,u_i-s. \]
Then $\phi'(t)=-2\sum_i(\lambda_iu_i+c_i)^2/(1+t\lambda_i)\leq0$, and if
$\phi(t^*)=0$, the point $x^0+Uu(t^*)$ is a point of $\Wc^H$ nearest to $y$.
\end{lem}
\begin{proof}
For $x=x^0+\hx$ with $\hx\perp z$,
$\|x-y\|^2=\|\hx-Py\|^2+\|x^0-(I-P)y\|^2$, and the second term does not
depend on $x$. For $t\in J$ the function $L(\hx)=\|\hx-Py\|^2+t(q(\hx)-s)$
is convex on $z^\perp$, with Hessian $2(I+t\hH)$, and its minimizer solves
$(I+t\hH)\hx=Py-t\hb$, i.e. $\hx=Uu(t)$, where $q(\hx)-s=\phi(t)$. If
$\phi(t^*)=0$, every $\hx'$ with $q(\hx')=s$ has
$\|\hx'-Py\|^2=L(\hx')\geq L(Uu(t^*))=\|Uu(t^*)-Py\|^2$. The derivative
follows from $u_i'=-(\lambda_iu_i+c_i)/(1+t\lambda_i)$.
\end{proof}

At the ends of $J$, $\phi$ tends to $+\infty$ and $-\infty$ unless
$\beta_i-tc_i$ vanishes there on the extreme eigenvectors, so in general it
has exactly one root. QUALEX-MS computes it with one eigendecomposition of
$\hH$ and a safeguarded Newton iteration per point
(\texttt{wrapper\_surface()} in \texttt{lib/wrapper.cc}), and its
Douglas--Rachford stage takes such surfaces as further sets: each start keeps
one copy on the sphere and one on the surface, and the orthant or the
2-clause projection acts on the average of their reflections, which is
Douglas--Rachford between the product of the sphere and the surface and the
diagonal of the third set (without the surface, the plain iteration). With
$H=H_0$ (\texttt{sat01qms -W}), on the instances of the previous subsection,
the orthant now finds the solution of 14471 and reaches 513 of $m=514$ on
60491 and 905 of 912 on 1000001, against 364, 505 and 895 without the
surface; with the 2-clause projection it solves the same seven instances as
before. The surface without nonnegativity (\texttt{sat01qms -N}) finds
nothing that the run without a Douglas--Rachford stage does not, and the
iteration converges poorly between the two quadrics: on the factoring
instances 1 to 19 of the 296 to 942 starts meet the stopping test
($10^{-9}$) within 300 iterations, and on the Hamilton cycle instances all
of them do, after 123 to 190 iterations on average (QUALEX-MS commits
892a830 and 02a809e, \texttt{bench/runs/sat01/qms.tsv}). Douglas--Rachford is
the wrong tool for that intersection, which is explicit:

\begin{prop}
\label{prop:cone}
Let $\Sigma=\{\bar x+Qp:\ p\in\R^k,\ \|p\|=\rho\}$ with $Q^TQ=I$ and
$Q^T\bar x=0$, let $H$ be symmetric, and write, for $x=\bar x+Qp$,
\[ x^THx-1=p^TMp+2g^Tp+h, \qquad M=Q^THQ,\ \ g=Q^TH\bar x,\ \
   h=\bar x^TH\bar x-1. \]
Let $M=V{\rm diag}(\kappa_i)V^T$, $\gamma=V^Tg$, and let $\sigma$, not an
eigenvalue of $M$, solve
\begin{equation}
\label{eq:conesigma}
\sum_i\frac{\gamma_i^2}{\kappa_i-\sigma}=h+\sigma\rho^2.
\end{equation}
Let $v_i=\gamma_i/(\sigma-\kappa_i)$ and $I_\pm=\{i:\ \pm(\kappa_i-\sigma)>0\}$,
and for unit vectors $\xi\in\R^{I_+}$ and $\eta\in\R^{I_-}$ let
$w_i=\xi_i/\sqrt{\kappa_i-\sigma}$ on $I_+$ and
$w_i=\eta_i/\sqrt{\sigma-\kappa_i}$ on $I_-$. Then the points of $\Sigma$ with
$x^THx=1$ are the points $\bar x+QV(v+\nu w)$ for all such $\xi,\eta$ and all
real roots $\nu$ of
\[ \|w\|^2\nu^2+2\,v^Tw\,\nu+\|v\|^2-\rho^2=0, \]
together with $\bar x+QVv$ if $\|v\|=\rho$. If $\|v\|<\rho$, each of these
equations has exactly one positive root $\nu_+$, and
$(\xi,\eta)\mapsto\bar x+QV(v+\nu_+w)$ maps the product of the unit spheres
of $\R^{I_+}$ and $\R^{I_-}$ one-to-one onto $\Sigma\cap\{x^THx=1\}$.
Equation (\ref{eq:conesigma}) has a root between any two consecutive
distinct eigenvalues of $M$ on whose eigenspaces $g$ does not vanish.
\end{prop}
\begin{proof}
On $\Sigma$, $\|p\|^2=\rho^2$, so for $u=V^Tp$, $K={\rm diag}(\kappa_i)$ and
every $\sigma$,
\[ x^THx-1=u^T(K-\sigma I)u+2\gamma^Tu+h+\sigma\rho^2
   =(u-v)^T(K-\sigma I)(u-v)-v^T(K-\sigma I)v+h+\sigma\rho^2, \]
since $(K-\sigma I)v=-\gamma$. As
$v^T(K-\sigma I)v=\sum_i\gamma_i^2/(\kappa_i-\sigma)$, (\ref{eq:conesigma})
makes the last two terms cancel, and on $\Sigma$
\[ x^THx-1=\sum_{i\in I_+}(\kappa_i-\sigma)(u-v)_i^2
   -\sum_{i\in I_-}(\sigma-\kappa_i)(u-v)_i^2. \]
This vanishes at $u=v$; for $u\neq v$ the two sums are not both zero, since
every index lies in $I_+\cup I_-$, so it vanishes iff they are equal to some
$\nu^2>0$, i.e. iff $u-v=\nu w$ with
$\xi_i=\sqrt{\kappa_i-\sigma}\,(u-v)_i/\nu$ and
$\eta_i=\sqrt{\sigma-\kappa_i}\,(u-v)_i/\nu$ unit vectors. Conversely each
such $w$ makes both sums equal to $1$, so $v+\nu w$ lies on the cone for
every $\nu$, and it lies on $\Sigma$ iff $\|v+\nu w\|=\rho$. If $\|v\|<\rho$,
the constant term of that quadratic is negative and its roots have opposite
signs. Distinct $(\xi,\eta)$ give distinct rays from $v$, because the
normalization fixes $w$ within its ray, and every point $u$ of the
intersection lies on the ray from $v$ through it ($u\neq v$, as
$\|v\|<\rho=\|u\|$), so the map is one-to-one and onto. Finally the left
side of (\ref{eq:conesigma}) minus its right side is continuous between two
consecutive distinct eigenvalues $\kappa_a<\kappa_b$, and tends to
$-\infty$ at $\kappa_a$ and to $+\infty$ at $\kappa_b$ when $g$ has a
component in both eigenspaces.
\end{proof}

The intersection thus has dimension $|I_+|+|I_-|-2=k-2$, and for
$\Sigma=\Sigma_A$ (with $\bar x$ its centre, $Q$ an orthonormal basis of
$\{x\in L_A:\ z^Tx=0\}$ and $\rho^2=1/m-\|\bar x\|^2$) and $H=H_0$ it costs
one eigendecomposition of order $k$ and one scalar equation.
\texttt{bench/cone.py} checks it (\texttt{bench/runs/sat01/cone-2026-10-10.txt}).
On 14471, where $k=148$, (\ref{eq:conesigma}) has a root in 145 of the 147
gaps of the spectrum of $M$. On the five instances checked (14471, 32399,
60491, the dodecahedron and $GP(13,2)$), the points of
Proposition~\ref{prop:cone} satisfy both equations to within
$2\cdot10^{-13}$, and the solution found by exact cover search lies on the
cone with its own $(\xi,\eta)$. But on all five every root puts the vertex
$v$ outside the ball, at $2.8\rho$ to $7.4\rho$ for the nearest, so the
generators that reach the sphere form a thin set: none of $10^5$ random ones
did, while those near a solution's do. And leaving the solution along the
intersection brings negative entries at once. Perturbing its $(\xi,\eta)$ by
$10^{-3}$ puts a median 3 to 5\% of the $\ell_1$ mass of the point on
negative entries on the factoring instances, and 16 to 19\% on the Hamilton
cycle ones.

So the two quadratic conditions cost an eigendecomposition and a root of
(\ref{eq:conesigma}), and by Proposition~\ref{prop:threesets} the solutions
are exactly the parameters $(\xi,\eta)$ whose point is nonnegative. All of
the hardness is in that last condition: when every contradiction lies in an
equation the surface removes no exact cover, and deciding whether the
intersection has a nonnegative point is the exact cover problem
(Corollary~\ref{cor:exactcover}).

\section{Computations}
\label{sec:comp}

\subsection{Tools}

The maximum cliques were listed with an unpublished maximal clique
enumerator (CLEN). The system (\ref{eq:system}) with a common level was set
up with repeated equations removed and tested for a solution at level $1$
(least squares, confirmed by linear programming); when solvable, the
spectrum was examined for the solution nearest the standard wrapper (least
change of the defects, level free), and nonnegative solutions were explored
by linear programming (largest smallest defect, and the defects that vanish
in every nonnegative solution). $\vartheta$ was computed with CSDP
\cite{Bo99} from whichever of two equivalent programs has fewer
constraints: (\ref{eq:thetadual}), one constraint per non-edge, or
$\min t$ subject to $tI-H\succeq0$, $H\in\Hc(G,w)$, i.e. a positive
semidefinite $Y=tI-H$ with $y_{ii}=t-1$ and $y_{ij}=-1$ on $E$, one
constraint per edge. The Lov\'asz-optimal $H$ was recovered from the solution.
The scripts are \texttt{bench/anchor\_all.py} and \texttt{bench/theta.py} of
the QUALEX-MS repository (commit c9af3f8), and their outputs are
\texttt{bench/runs/anchor-all-2026-09-26.txt} and
\texttt{bench/runs/theta/theta.tsv}.

The following checks agree with the statements above:
Theorem~\ref{thm:level} to $4\cdot10^{-13}$ on 127 pairs of maximal cliques
of weighted $G(40,0.5)$ that admit nonzero levels; the criterion of
Theorem~\ref{thm:two} on all 150 pairs tried (23 of which admit only level
$0$); Proposition~\ref{prop:curv} to $10^{-15}$ on 1997 pairs of cliques;
Theorem~\ref{thm:collision} with multiplicity equal to $d$ and $|c|\leq
10^{-16}$ on the eigenspace; Theorem~\ref{thm:holes} and the even hole
statement for $C_5,\ldots,C_{10}$ and $\bar C_7,\bar C_9,\bar C_{11}$; for
$C_5$, over a grid of the free parameters of one anchored edge,
$\lmax(\hH)-\mu\geq0.138$ (Theorem~\ref{thm:global}, $\vartheta=\sqrt5>2$);
the colouring wrapper on random interval graphs; and
Theorem~\ref{thm:theta} and Corollary~\ref{cor:hoffman} on the graphs of
Table~\ref{tab:bounds}. CSDP was validated on \gr{hamming6-2} (32),
\gr{johnson8-4-4} (14), \gr{hamming6-4} ($16/3$) and \gr{MANN_a9}
($17.47503$), with both programs and with CSDP's own \texttt{theta} program.

\subsection{Anchoring all maximum cliques of DIMACS graphs}

Table~\ref{tab:dimacs} covers every DIMACS graph with $n\leq1100$ and
between 2 and 20000 maximum cliques. For 11 of them the common level is
forced to $0$; by Corollary~\ref{cor:certificate} all have $\vartheta>
\omega$. For the other 20 the level is free. On 6 of these the anchoring
nearest the standard wrapper already puts $\mu$ at the top of the spectrum
of $\hH$, which by Theorem~\ref{thm:global} certifies $\vartheta=\omega$;
for five of them the standard wrapper itself anchors all maximum cliques,
because every outside vertex has the same number of non-neighbours in each.
(For \gr{johnson8-4-4} the maximum cliques are the 30 Steiner systems
$S(3,4,8)$, and a 4-set outside such a system meets exactly 4 of its 14
blocks in three points.) On the remaining 14 the nearest anchoring leaves
between 2 and 232 eigenvalues above $\mu$, and CSDP decides:
$\vartheta=\omega$ on 8 of them and $\vartheta>\omega$ on 6. Wherever the
level is free, the multiplicity of $\mu$ in $\hH$ equals $d$ exactly, as
Theorem~\ref{thm:collision} predicts as a lower bound, with the single
exception of \gr{hamming8-4} (127 against 71).

\begin{table}[p]
\begin{center}
\small
\begin{tabular}{|l|r|r|r|r|l|r|l|}
\hline
graph & $n$ & $\omega$ & $k$ & $d$ & common level & above $\mu$ & $\vartheta(\bar G)$ \\
\hline
\gr{p_hat300-1}   & 300  & 8   & 13   & 12  & forced to 0 & & $>\omega$ \\
\gr{sanr200_0.7}  & 200  & 18  & 13   & 10  & forced to 0 & & $>\omega$ \\
\gr{p_hat300-2}   & 300  & 25  & 52   & 17  & forced to 0 & & $>\omega$ \\
\gr{DSJC500.5}    & 500  & 13  & 56   & 55  & forced to 0 & & $>\omega$ \\
\gr{p_hat500-1}   & 500  & 9   & 78   & 74  & forced to 0 & & $>\omega$ \\
\gr{p_hat700-2}   & 700  & 44  & 138  & 23  & forced to 0 & & $>\omega$ \\
\gr{hamming6-4}   & 64   & 4   & 240  & 51  & forced to 0 & & 5.33333 \\
\gr{p_hat1000-1}  & 1000 & 10  & 276  & 260 & forced to 0 & & $>\omega$ \\
\gr{C125.9}       & 125  & 34  & 924  & 24  & forced to 0 & & $>\omega$ \\
\gr{keller4}      & 171  & 11  & 2304 & 133 & forced to 0 & & $>\omega$ \\
\gr{MANN_a9}      & 45   & 16  & 9540 & 44  & forced to 0 & & 17.47503 \\
\hline
\gr{hamming6-2}   & 64   & 32  & 2    & 1   & free, standard & 0 & 32 \\
\gr{hamming8-2}   & 256  & 128 & 2    & 1   & free, standard & 0 & $=\omega$ \\
\gr{hamming10-2}  & 1024 & 512 & 2    & 1   & free, standard & 0 & $=\omega$ \\
\gr{johnson8-4-4} & 70   & 14  & 30   & 14  & free, standard & 0 & 14 \\
\gr{johnson8-2-4} & 28   & 4   & 105  & 20  & free, standard & 0 & $=\omega$ \\
\gr{hamming8-4}   & 256  & 16  & 480  & 71  & free & 0 & $=\omega$ \\
\hline
\gr{san200_0.7_2}   & 200 & 18  & 2  & 1  & free & 43  & 18 \\
\gr{gen200_p0.9_55} & 200 & 55  & 4  & 2  & free & 13  & 55 \\
\gr{gen200_p0.9_44} & 200 & 44  & 4  & 2  & free & 33  & 44 \\
\gr{c-fat200-1}     & 200 & 12  & 14 & 13 & free & 8   & 12 \\
\gr{c-fat500-1}     & 500 & 14  & 19 & 18 & free & 19  & 14 \\
\gr{c-fat500-2}     & 500 & 26  & 19 & 18 & free & 7   & 26 \\
\gr{c-fat500-5}     & 500 & 64  & 3  & 2  & free & 4   & 64 \\
\gr{c-fat500-10}    & 500 & 126 & 3  & 2  & free & 2   & 126 \\
\gr{brock200_1}     & 200 & 21  & 2  & 1  & free & 49  & 27.45664 \\
\gr{p_hat700-1}     & 700 & 11  & 2  & 1  & free & 232 & 15.11986 \\
\gr{c-fat200-5}     & 200 & 58  & 3  & 2  & free & 2   & 60.34528 \\
\gr{sanr400_0.5}    & 400 & 13  & 4  & 3  & free & 143 & 20.31766 \\
\gr{p_hat300-3}     & 300 & 36  & 10 & 5  & free & 90  & 41.16993 \\
\gr{p_hat500-2}     & 500 & 36  & 14 & 5  & free & 151 & 38.97466 \\
\hline
\end{tabular}
\end{center}
\caption{\label{tab:dimacs} DIMACS graphs with $2\leq k\leq20000$ maximum
cliques and $n\leq1100$. $d$ is the dimension of the affine hull of their
indicators; ``standard'': the standard wrapper already anchors them all;
``above $\mu$'': eigenvalues of $\hH$ above the common level for the
anchoring nearest the standard wrapper. $\vartheta$: computed with CSDP
where a value is given; $>\omega$ by Corollary~\ref{cor:certificate};
$=\omega$ by Theorem~\ref{thm:global} (and, for the hamming graphs of
distance 2, which are complements of bipartite graphs, by perfection).}
\end{table}

Two features of the solutions are worth recording. First,
Proposition~\ref{prop:blind} is not an exception. Only 7 of the 20 graphs,
the six symmetric ones and \gr{brock200_1}, admit a proper anchoring of all
their maximum cliques; on the other 13, between 2.6\% (\gr{san200_0.7_2})
and 23\% (\gr{c-fat200-5}) of the non-edges involved are invisible in every
nonnegative solution. (The solutions $(\tau,\mu)$ form a cone, so a single
linear program finds these non-edges: maximize $\sum_k\min(\tau_k,1)$
subject to (\ref{eq:system}), $\tau\geq0$ and $\mu\geq1$; exactly the
defects that can be positive reach $1$.) Second, in the graphs with the
level forced to $0$ the collapse is local and has small certificates, of the
kind of Example~\ref{ex:cascade}.

\subsection{Random graphs}

On the 35 unweighted random graphs with proven $\omega$ ($n=200$ with
$p\leq0.8$ and $n=300$ with $p\leq0.7$), 5 have a unique maximum clique. Of
the 30 others the common level is forced to $0$ on 17 (with $k$ between 3
and 797) and free on 13 (with $k$ between 2 and 15); where free, $\mu$ had
55 to 99 eigenvalues of $\hH$ above it and the multiplicity of $\mu$ was $d$
exactly. One of the 13 admitted no nonnegative solution, 5 admitted proper
ones, and on the other 7 between 2.1\% and 9.4\% of the non-edges involved
are invisible in every nonnegative solution. Weighted random graphs almost always have a unique maximum weight
clique, for which anchoring is the single clique construction of
Section~\ref{sec:one}.

\subsection{The Lov\'asz number of the 14 undecided graphs}

Table~\ref{tab:theta} lists the CSDP results for the 14 graphs on which the
nearest anchoring stayed below the top. In all eight graphs with
$\vartheta=\omega$ the optimal wrapper anchors every maximum clique at $C=0$
(to $10^{-9}$ to $10^{-4}$), as Theorem~\ref{thm:global} requires. The
multiplicity of $\lmax(H)$ agrees with the rank of the Lov\'asz optimum found
by CSDP (strict complementarity) and, but for one graph, with
$\dim\spn\{z_Q\}$, the lower bound of Corollary~\ref{cor:hard}. The
exception is \gr{gen200_p0.9_44}: its Lov\'asz optimum has rank 157 and a
positive diagonal on all 200 vertices, so by (\ref{eq:compl}) every
Lov\'asz-optimal wrapper has a top eigenvalue of multiplicity at least 157,
though its four maximum cliques span only 3 dimensions. The optimal face of
(\ref{eq:thetadual}) contains ``vector cliques'' that are not cliques. When
$\vartheta>\omega$ the top multiplicities are large and unrelated to the
cliques; on the largest graphs the top of the spectrum decays gradually and
the multiplicity is only defined up to the solver's accuracy.

\begin{table}[ht]
\begin{center}
\small
\begin{tabular}{|l|r|r|r|r|r|}
\hline
graph & $\omega$ & $\vartheta(\bar G)$ & $\dim\spn\{z_Q\}$ & rank of $X^*$ & mult.\ of $\lmax(H)$ \\
\hline
\gr{c-fat200-1}     & 12  & 12.000000  & 14 & 14  & 14 \\
\gr{c-fat500-1}     & 14  & 14.000000  & 19 & 19  & 19 \\
\gr{c-fat500-2}     & 26  & 26.000000  & 19 & 19  & 19 \\
\gr{c-fat500-5}     & 64  & 64.000001  & 3  & 3   & 3 \\
\gr{c-fat500-10}    & 126 & 126.000000 & 3  & 3   & 3 \\
\gr{gen200_p0.9_55} & 55  & 55.000000  & 3  & 3   & 3 \\
\gr{gen200_p0.9_44} & 44  & 44.000000  & 3  & 157 & 157 \\
\gr{san200_0.7_2}   & 18  & 18.000000  & 2  & 2   & 2 \\
\hline
\gr{brock200_1}     & 21  & 27.456641  & 2  & 63  & 63 \\
\gr{c-fat200-5}     & 58  & 60.345275  & 3  & 5   & 5 \\
\gr{p_hat300-3}     & 36  & 41.169930  & 6  & 101 & 101 \\
\gr{sanr400_0.5}    & 13  & 20.317656  & 4  & $\approx200$ & $\approx200$ \\
\gr{p_hat500-2}     & 36  & 38.974658  & 6  & $\sim150$ & $\sim150$ \\
\gr{p_hat700-1}     & 11  & 15.119858  & 2  & $\sim450$ & $\sim450$ \\
\hline
\end{tabular}
\end{center}
\caption{\label{tab:theta} The Lov\'asz number of the complement (CSDP,
relative duality gaps below $10^{-8}$), the dimension spanned by the maximum
cliques, the rank of the Lov\'asz optimum $X^*$ and the multiplicity of the
top eigenvalue of the optimal wrapper, both read at the largest relative
gap in the spectrum. For \gr{sanr400_0.5} that gap lies between 199 and 201.
For the last two graphs there is no clear gap: counting eigenvalues within
relative distances $10^{-6}$ to $10^{-4}$ gives 74 to 165 (\gr{p_hat500-2})
and 428 to 486 (\gr{p_hat700-1}). The largest problems had 61800
constraints and took 42 and 47 minutes.}
\end{table}

\subsection{Relaxed bounds of the standard and the optimal wrapper}

Table~\ref{tab:bounds} compares the bound $1/v(H)$ of
Theorem~\ref{thm:theta} for the standard wrapper and for the optimal
wrapper found by CSDP; $v(H)$ was computed from the secular equation as the
stationary point with the largest multiplier (Lemma~\ref{lem:global}),
including the hard case. For the optimal wrappers $1/v(H)=\vartheta$ to the
accuracy of the solver (relative error below $3\cdot10^{-6}$); on the three
regular graphs $1/v(H_0)$ equals Hoffman's bound, as
Corollary~\ref{cor:hoffman} states; elsewhere the standard wrapper's bound
can be very weak (\gr{san200_0.7_2}: 115.8 against $\vartheta=18$).

\begin{table}[ht]
\begin{center}
\small
\begin{tabular}{|l|r|r|r|r|}
\hline
graph & $\omega$ & $\vartheta(\bar G)$ & $1/v(H_0)$ & Hoffman \\
\hline
\gr{hamming6-2}     & 32 & 32.000 & 32.000  & 32.000 \\
\gr{johnson8-4-4}   & 14 & 14.000 & 14.000  & 14.000 \\
\gr{hamming6-4}     & 4  & 5.333  & 13.538  & 13.538 \\
\gr{MANN_a9}        & 16 & 17.475 & 19.707  & \\
\gr{c-fat200-1}     & 12 & 12.000 & 17.430  & \\
\gr{gen200_p0.9_55} & 55 & 55.000 & 75.624  & \\
\gr{gen200_p0.9_44} & 44 & 44.000 & 72.547  & \\
\gr{san200_0.7_2}   & 18 & 18.000 & 115.759 & \\
\gr{brock200_1}     & 21 & 27.457 & 43.287  & \\
\gr{c-fat200-5}     & 58 & 60.345 & 72.905  & \\
\gr{p_hat300-3}     & 36 & 41.170 & 88.374  & \\
\hline
\end{tabular}
\end{center}
\caption{\label{tab:bounds} The relaxed bound $1/v(H)$ for the standard
wrapper against $\vartheta$, which the optimal wrapper attains
(Theorem~\ref{thm:theta}), and Hoffman's bound for the regular graphs.}
\end{table}

\subsection{Quasigroups with holes}
\label{sec:qwh}

Proposition~\ref{prop:qcp} places the least top multiplicity of a
Lov\'asz-optimal wrapper of a completable QCP graph between the span of the
completions and $\dim L$. To see where it falls, we generated
quasigroup-with-holes (QWH) instances \cite{AGKS00}: a random latin square of
order $n$ from the Jacobson--Matthews chain \cite{JM96}, with $h$ cells
emptied uniformly at random, for $n=10,15,20$ and hole fractions from 30 to
70\% (four instances per cell, three for $n=20$). The range covers the
passage from a single completion to tens of thousands, where QWH instances
are hardest for search \cite{AGKS00}. The completions were listed by exact
cover search; where there were more than 20000 of them, the span was
measured on the first 20000 together with about 2000 sampled by randomized
descents, restarted after a node limit because such searches have
heavy-tailed run times. $\vartheta$ and a Lov\'asz optimum were computed with CSDP.
Table~\ref{tab:qwh} summarizes the 68 instances; \texttt{bench/qwh.py}
reproduces them.

Every instance confirms Proposition~\ref{prop:qcp}: $\vartheta(G)=h$, the
standard wrapper in the gauge $J-\frac h3A_G$ has the top eigenvalue $h$
with multiplicity exactly $\dim L$, and every completion is anchored at
$C=0$ to $10^{-13}$.

The Lov\'asz optimum spans exactly the completions whenever there are few of
them: its rank equals their span on 29 of the 41 instances listed in full,
among them all nine with a unique completion, where the optimum found by
CSDP is the projector of that completion. It does so again when there are
very many: on the 27 instances whose listing was capped, the listed and
sampled completions span as much as the optimum on 26 of them and one
dimension less on the other, and at 70\% holes both fill all of $L$. In between, with
tens to hundreds of completions, the optimum exceeds their span: 64
dimensions against 34, 40 against 18 and 75 against 60 on three instances
with $n=15$ and 45\% holes, 15 against 7 on one with $n=10$. Every
Lov\'asz-optimal wrapper of those graphs has top eigendirections that no
completion spans. On a few of them complementarity is not strict: some
directions carry both an eigenvalue of the optimum and a gap of the wrapper
of order $\sqrt\mu$, and the rank is only bracketed (the parenthesized
ranges of Table~\ref{tab:qwh}).

The natural wrappers are far more degenerate. The standard wrapper's top
eigenspace $L$ had dimension 111 on an instance whose two completions span 2
dimensions ($n=20$, 35\% holes), and 194 against a span of 4 on another
($n=20$, 40\% holes); the colouring wrapper of the cells had 327 and 453.
So on such instances only a handful of the degenerate directions of the
standard QUALEX-MS wrapper carry information on the completions, as
\cite{BP03} anticipates. The least degenerate Lov\'asz-optimal wrapper can be
much sharper, but it too exceeds the completions in the middle of the
transition. These instances are small: all nine with a unique completion
were solved by propagation alone, so they say nothing about hard instances
with a unique completion.

\begin{table}[ht]
\begin{center}
\small
\setlength\tabcolsep{4pt}
\begin{tabular}{|r|r|r|r|l|r|l|r|r|}
\hline
$n$ & holes & inst. & $|V|$ & completions & span & rank of $X^*$ & $\dim L$ & $|V|-h+1$ \\
\hline
10 & 30\% & 4 & 50--59 & 1--2 & 1--2 & 1--2 & 1--3 & 21--30 \\
10 & 35\% & 4 & 73--81 & 1--4 & 1--3 & 1--3 & 4--11 & 39--47 \\
10 & 40\% & 4 & 87--92 & 1--2 & 1--2 & 1--2 & 6--9 & 48--53 \\
10 & 45\% & 4 & 126--143 & 3--15 & 3--12 & 3--13 & 22--39 & 82--99 \\
10 & 50\% & 4 & 154--164 & 4--11 & 3--8 & 3--15 (3--17) & 34--44 & 105--115 \\
10 & 60\% & 4 & 241--261 & $>$20000 & 88--101 & 88--101 & 91--111 & 182--202 \\
10 & 70\% & 4 & 369--379 & $>$20000 & 189--199 & 189--199 & 189--199 & 300--310 \\
15 & 30\% & 4 & 145--170 & 1--2 & 1--2 & 1--2 & 5--18 & 78--103 \\
15 & 35\% & 4 & 192--216 & 1--4 & 1--4 & 1--4 & 12--36 & 114--138 \\
15 & 40\% & 4 & 277--312 & 14--38 & 10--15 & 11--19 (14--23) & 56--87 & 188--223 \\
15 & 45\% & 4 & 369--381 & 49--5029 & 18--104 & 40--105 (41--105) & 111--125 & 269--281 \\
15 & 50\% & 4 & 485--520 & $>$20000 & 176--192 & 176--192 & 194--230 & 374--409 \\
15 & 60\% & 4 & 780--805 & $>$20000 & 416--437 & 416--437 & 420--445 & 646--671 \\
15 & 70\% & 4 & 1216--1260 & $>$20000 & 787--831 & 787--831 & 787--831 & 1059--1103 \\
20 & 35\% & 3 & 452--466 & 2--9 & 2--8 & 2--9 & 102--111 & 313--327 \\
20 & 40\% & 3 & 612--684 & 8--7132, $>$20000 & 4--218 & 4--218 & 194--266 & 453--525 \\
20 & 45\% & 3 & 837--868 & $>$20000 & 329--360 & 329--360 & 358--388 & 658--689 \\
20 & 50\% & 3 & 1120--1140 & $>$20000 & 533--585 & 533--585 & 580--600 & 921--941 \\
\hline
\end{tabular}
\end{center}
\caption{\label{tab:qwh} QWH instances, ranges over the instances of a cell.
$|V|$: vertices of the QCP graph; completions: listed by exact cover,
$>$20000 where the listing was capped (the span then also covers about 2000
sampled completions); span: dimension spanned by the completion indicators;
rank of the Lov\'asz optimum $X^*$ found by CSDP, counting eigenvalues above
$10^{-4}$ of the largest (in parentheses above $10^{-6}$, where that
differs); $\dim L$: top multiplicity of the standard wrapper; $|V|-h+1$: top
multiplicity of the colouring wrapper of the cells.}
\end{table}

\subsection{SAT01 instances}
\label{sec:sat01comp}

Two reductions of \cite{Bu07} supplied the instances, built with Busygin's
converters. A factoring instance encodes the long multiplication of two
factors with product $N$, with the bitwise sums of
\cite[Section~4.9]{Bu07}; trivial factorizations are excluded, so a prime
gives an unsatisfiable instance, and the products of two primes tried had
one solution, or two when both factors have the same bit length. A Hamilton
cycle instance of a graph with $n$ vertices has the variables $x_{ij}$ (``the
$i$-th vertex of the cycle is vertex $j$''), $2n$ equations that make
$(x_{ij})$ a permutation matrix, and a contradiction between $x_{ij}$ and
$x_{i+1,k}$ whenever $j$ and $k$ are not adjacent \cite[Section~4.2]{Bu07};
one variable is fixed to remove the rotations, and once it is eliminated
every weight is $2$ and $m=2(n-1)$. The graphs were generalized Petersen
graphs $GP(n,k)$, among them the Petersen graph $GP(5,2)$ and the
dodecahedron $GP(10,2)$, the flower snarks $J_5$ and $J_7$, the Coxeter
graph, and random cubic graphs; the solver confirmed that the Petersen
graph, $GP(11,2)$, $J_5$, $J_7$ and the Coxeter graph have no Hamilton cycle.

Each instance was converted after SAT01 propagation of two strengths: {\em
light}, which only eliminates the variables that the encoding fixes, and
{\em full}, unit propagation, clique analysis and contradiction analysis to a
fixed point, as in Busygin's solver. Full propagation refutes all 1728 primes
between $2^{10}$ and $2^{14}$ but 14887, four of the eight largest primes
below $2^{15}$, and none of the eight largest below $2^{16}$ or below
$2^{17}$. It solves the 2957 odd products of two primes in the same range or
reduces them to one equation in two variables, except 14471, 14631, 15591
and 16041. Bipartite graphs with sides of different sizes, which have no
Hamilton cycle by parity, are refuted already by light propagation, whose
clique analysis carries the parity of the positions from the fixed vertex.

$\vartheta$ was computed with CSDP from (\ref{eq:thetadual}), one constraint
per contradiction. A value below $m$ is certified by $\lmax$ of the wrapper
recovered from the dual solution, which bounds $\vartheta$ from above by
(\ref{eq:thetaprimal}) whatever the solver's accuracy. The rank of a
Lov\'asz optimum counts its eigenvalues above $10^{-5}$ of the largest, or
those before a drop by more than three orders of magnitude where there is one
(only $179\cdot181$ has one above $10^{-5}$), the multiplicity of $\lmax(H)$
counts the eigenvalues within $10^{-5}\lmax(H)$ of it, and the solutions were
listed by exact cover search. The script is
\texttt{bench/sat01.py} of the QUALEX-MS repository (commits b51234b and
33020d1), its output is in \texttt{bench/runs/sat01/}, and the largest
problems, with 55000 constraints, took 27 minutes.

\begin{table}[ht]
\begin{center}
\small
\setlength\tabcolsep{3pt}
\begin{tabular}{|l|r|r|r|l|r|r|r|r|l|r|}
\hline
 & & \multicolumn{3}{c|}{light propagation} & \multicolumn{6}{c|}{full propagation} \\
instance & guesses & $n$ & $m$ & $m-\vartheta$ & $n$ & $m$ & contr. & blind & $m-\vartheta$ & rank $X^*$ \\
\hline
$101$   & 0 & 145 & 106 & $0.103^*$ & \multicolumn{6}{c|}{refuted by propagation} \\
$1009$  & 0 & 316 & 217 & 0 & \multicolumn{6}{c|}{refuted by propagation} \\
$14887$ & 1 & 716 & 478 & 0 & 510 & 365 & 22926 & 0.93 & $0.0638^*$ & 94 \\
$32717$ & 1 & 868 & 575 & 0 & 568 & 404 & 29157 & 0.94 & $0.0345^*$ & 108 \\
$32719$ & 1 & 866 & 575 & 0 & 615 & 425 & 30173 & 0.93 & $8.9\cdot10^{-4}\,{}^*$ & 101 \\
$32687$ & 1 & 866 & 575 & 0 & 638 & 435 & 34862 & 0.94 & $7.0\cdot10^{-5}\,{}^*$ & 106 \\
$32713$ & 1 & 871 & 575 & 0 & 612 & 419 & 29772 & 0.93 & 0 & 161 \\
$65521$ & 1 & 955 & 631 & 0 & 713 & 482 & 41260 & 0.94 & 0 & 198 \\
$32749,32707,32693,32653$ & 0 & 865--872 & 575 & 0 & \multicolumn{6}{c|}{refuted by propagation} \\
\hline
Petersen   & 0  & 63 & 18 & $0.235^*$ & \multicolumn{6}{c|}{refuted by propagation} \\
$J_5$      & 2  &    &    &   & 265 & 38 & 14796 & 0.74 & 0 & 211 \\
$GP(11,2)$ & 5  &    &    &   & 329 & 42 & 19316 & 0.73 & 0 & 269 \\
Coxeter    & 83 &    &    &   & 585 & 54 & 50832 & 0.74 & 0 & 509 \\
$J_7$      & 30 &    &    &   & 543 & 54 & 55123 & 0.79 & 0 & 465 \\
\hline
\end{tabular}
\end{center}
\caption{\label{tab:sat01unsat} Unsatisfiable SAT01 instances: primes
(factoring) and graphs without a Hamilton cycle. guesses: made by the SAT01
solver, 0 when propagation refutes the instance. After light and after full
propagation: $n$ variables, $m$ equations, the contradictions, the share of
them that share no equation (blind to $H_A$), $m-\vartheta$, and the rank of
the Lov\'asz optimum $X^*$. A $0$ means $|m-\vartheta|\leq5\cdot10^{-6}$ with
the dual wrapper's $\lmax$ above $m$; $^*$: certified by a dual wrapper with
$\lmax<m$.}
\end{table}

\begin{table}[ht]
\begin{center}
\small
\begin{tabular}{|l|r|r|r|r|r|r|r|}
\hline
instance & $n$ & $m$ & solutions & span & rank $X^*$ & mult.\ $\lmax(H)$ & $\dim L_A$ \\
\hline
$143=11\cdot13$ (light) & 180 & 132 & 2  & 2  & 2   & 2   & 49 \\
$14471=29\cdot499$      & 511 & 367 & 1  & 1  & 1   & 1   & 149 \\
$14631=3\cdot4877$      & 514 & 370 & 1  & 1  & 1   & 2   & 150 \\
$15591=3\cdot5197$      & 501 & 359 & 1  & 1  & 1   & 1   & 145 \\
$16041=3\cdot5347$      & 506 & 362 & 1  & 1  & 1   & 1   & 146 \\
$32399=179\cdot181$     & 636 & 441 & 2  & 2  & 2   & 18  & 198 \\
$60491=241\cdot251$     & 740 & 514 & 2  & 2  & 205 & 211 & 236 \\
\hline
cubic, 16 vertices      & 57  & 30  & 12 & 12 & 12  & 12  & 29 \\
cubic, 20 vertices      & 195 & 38  & 30 & 30 & 152 & 152 & 160 \\
cubic, 24 vertices      & 403 & 46  & 48 & 36 & 335 & 335 & 359 \\
dodecahedron            & 263 & 38  & 60 & 60 & 213 & 213 & 227 \\
$GP(13,2)$              & 481 & 50  & 26 & 26 & 409 & 409 & 433 \\
\hline
\end{tabular}
\end{center}
\caption{\label{tab:sat01sat} Satisfiable SAT01 instances after full
propagation ($143$ after light propagation; full propagation leaves one
equation in two variables). solutions: listed by exact cover search, the
directed Hamilton cycles from the fixed vertex for the graphs; span: the
dimension of their indicators; rank of the Lov\'asz optimum $X^*$;
multiplicity of the top eigenvalue of the optimal wrapper; $\dim L_A$: that of
$H_A$. Every solution is anchored at $C=0$ in the optimal wrapper to a
relative $4\cdot10^{-5}$. For $179\cdot181$ the eigenvalues of $X^*$ fall by a
factor 5000 after the second and then decay smoothly from $5\cdot10^{-5}$ of the
largest, the $\sqrt\mu$ tail of a solution that is not strictly
complementary; counted above $10^{-5}$ they would give rank 5.}
\end{table}

Table~\ref{tab:sat01unsat} answers the question of \cite{Bu07} on these
instances. After full propagation $\vartheta$ refutes four of the six primes
that propagation leaves open, each certified by a wrapper with $\lmax<m$,
but the margin $m-\vartheta$ falls from $0.064$ (14887) to $7\cdot10^{-5}$
(32687). On 32713 and 65521 $\vartheta=m$: by Theorem~\ref{thm:vector} they
have vector solutions, of dimension about 160 and 200, and they are
unsatisfiable instances with $\alpha<\vartheta=m$. On the graphs of light
propagation $\vartheta=m$ for all ten primes from 14887 up, while it still
refutes 101, so each of the four refutations needs the contradictions that
contradiction analysis derives. Conversely, propagation alone refutes 32749,
32707, 32693 and 32653, whose light graphs have $\vartheta=m$. Neither method
contains the other, and propagation followed by $\vartheta$ refutes more than
either. The polyhedral shadow of Theorem~\ref{thm:vector}, a fractional
exact cover with $p_i+p_j\leq1$ on the contradictions, was feasible on all
seven instances tried, among them the three with $\vartheta<m$ (101 light,
14887 full, Petersen light), so these refutations are semidefinite and not
polyhedral. For perspective, the solver needs a single guess on each
surviving prime: the instances are easy, and $\vartheta$ helps at the edge of
what propagation does. The 17-bit primes keep 60000 to 63000 contradictions
after propagation, beyond what CSDP solves on the machine used.

On the Hamilton cycle instances $\vartheta$ refutes only the Petersen graph
($17.765<18$ on its light graph), which propagation refutes anyway. The four
non-Hamiltonian graphs that survive full propagation all have
$\vartheta=m$: they have vector Hamilton cycles, of dimension 211 to 509.

On the satisfiable instances (Table~\ref{tab:sat01sat}) the solutions are
the vector solutions of dimension one, and the question is how many others
the Lov\'asz optimum carries. The four 14-bit products of two primes have a
unique solution and a Lov\'asz optimum of rank one: the solution is their
only vector solution, and solving the semidefinite program factors them.
The optimum of $179\cdot181$ is still spanned by its two solutions, but those
of $241\cdot251$ lie in an optimum of rank 205. On the cubic graphs the
optimum spans exactly the 12 solutions of the smallest one (6 cycles, each in
two directions) and far more than the solutions elsewhere: 152 dimensions
against 30, 335 against 36, 213 against 60 on the dodecahedron and 409
against 26 on $GP(13,2)$, as on the QWH instances in the middle of the
transition (Section~\ref{sec:qwh}). The optimal wrapper's top multiplicity
equals the rank on all the graphs and exceeds it on three factoring
instances (2 against 1, 18 against 2, 211 against 205), where
complementarity is not strict at the solver's accuracy. $H_A$, which is
Lov\'asz-optimal on all of them by Proposition~\ref{prop:sat01}(c), is far
more degenerate on factoring, with $\dim L_A$ between 145 and 236 against
ranks between 1 and 205, but exceeds the optimum by only 8 to 24 dimensions
on the Hamilton cycle instances.

\section{Conclusions and open questions}

For one clique, anchoring is always possible and puts the clique on the
secular curve; its effect is local, with ascent towards heavier cliques in
proportion to the gain per exchanged weight. Used on the solver's own best
clique it is a reliable improvement step, used on a poor clique it is
harmful. For several cliques, the free entries are far fewer in effect than
in number: the levels are forced by the level identity, cliques of equal
weight collide into the degenerate case, and the remaining system may force
non-edges to be invisible or the level to vanish. Anchoring all maximum
cliques at the top of the spectrum is possible exactly when
$\vartheta=\omega$, and then it is what every Lov\'asz-optimal wrapper does;
this explains the degeneracy of the Lov\'asz-optimal wrappers observed in
QUALEX-MS, even though by Theorem~\ref{thm:theta} they are the best
wrappers for the relaxed program. On odd holes and odd antiholes no
nontrivial wrapper anchors all maximum cliques, and this
characterizes perfect graphs through the strong perfect graph theorem.

The quasigroup completion graphs of Busygin and Pasechnik \cite{BP03} show
why the hard case cannot be put to work in general. For a completable
instance the standard wrapper that QUALEX-MS builds is itself Lov\'asz-optimal,
its degenerate top eigenspace contains all the completions, and they are
exactly the nonnegative points of its sphere of global minimizers
(Proposition~\ref{prop:qcp}); finding one of them is QCP. On random
instances that eigenspace is often much larger than the span of the
completions, and even the least degenerate Lov\'asz-optimal wrapper exceeds
that span in the middle of the transition from one completion to many
(Section~\ref{sec:qwh}).

SAT01 \cite{Bu07} carries this over to every problem that reduces to it
(Section~\ref{sec:sat01}). The equation wrapper $H_A$ is Lov\'asz-optimal
exactly when $\vartheta=m$; then its degenerate top eigenspace holds every
vector solution, and the nonnegative points of its sphere of global
minimizers are the exact covers of the equations, blind to the other
contradictions. $\vartheta<m$ refutes an instance exactly when it has no
vector solution (Theorem~\ref{thm:vector}). On factoring and Hamilton cycle
instances vector solutions appear early (Section~\ref{sec:sat01comp}).
$\vartheta$ refutes small primes only on the graphs that SAT01 propagation
has enriched with derived contradictions, with a margin that shrinks fast,
and no non-Hamiltonian graph that survives propagation; propagation in turn
refutes primes on whose unpropagated graphs $\vartheta=m$.

The NP-completeness is therefore not hidden in the wrappers. For a proper
Lov\'asz-optimal wrapper the maximum cliques are exactly the nonnegative
points of the sphere of global minimizers (Proposition~\ref{prop:proper}),
and exact cover has such a wrapper, $H_A$, on every satisfiable instance,
without any semidefinite program. Deciding whether that sphere has a
nonnegative point is the exact cover problem itself
(Corollary~\ref{cor:exactcover}). The semidefinite program and linear algebra
produce the polytope of nonnegative points of the degenerate eigenspace,
and what remains NP-hard is finding a vertex of largest norm, which is the
step from a Lov\'asz optimum to a rank-one one. Contradictions outside the
equations add improperness (Proposition~\ref{prop:improper}), but that is
not what keeps the problem hard. Nor do the quadratic conditions. The
surface of the standard wrapper, which every solution lies on, cuts the
sphere down to the solutions together with nonnegativity
(Proposition~\ref{prop:threesets}), and its intersection with the sphere is
explicit (Proposition~\ref{prop:cone}). So everything but nonnegativity
costs linear algebra and one scalar equation. In QUALEX-MS the surface helps
the orthant (Section~\ref{sec:surface}), but without nonnegativity it finds
nothing, and the points of the intersection turn negative as soon as they
leave a solution.

Questions left open:
\begin{enumerate}
\item Characterize the graphs whose maximum cliques can all be anchored at a
nonzero level; the condition is not hereditary, and Theorem~\ref{thm:two}
settles only pairs.
\item Given that $\vartheta>\omega$, how few eigenvalues can lie above the
common level of an anchoring? This is an eigenvalue optimization over the
affine family of anchoring wrappers.
\item Is there a weighted analogue of Theorem~\ref{thm:perfect}, i.e. which
weighted graphs have, hereditarily, anchorings of all maximum weight
cliques?
\item What structure makes the Lov\'asz optimum of \gr{gen200_p0.9_44} full
support and rank 157, and is it typical of generated instances with a
planted clique?
\item On QCP graphs, when does the Lov\'asz optimum span more than the
completions? On the small instances of Section~\ref{sec:qwh} it did only
between the regimes of few and of very many completions, and none of them
was hard for search. (An earlier version of these notes asked whether the
hard case could be exploited by a search inside the degenerate eigenspace;
Sections~\ref{sec:qcp} and~\ref{sec:hardness} answer that it cannot in
general, unless P = NP.)
\item Which primes keep a vector solution after SAT01 propagation (32713 and
65521 did, while 32687 lost it by a margin of $7\cdot10^{-5}$), and does every
prime from some size on keep one?
\item Contradiction analysis and $\vartheta$ are incomparable. Declaring $x_j$
and $x_k$ contradictory when the instance with $x_j=x_k=1$, after
propagation, has $\vartheta$ below its number of equations is a polynomial
rule that subsumes contradiction analysis. How much does it add on the
instances of Section~\ref{sec:sat01comp}, and does iterating it refute
32713?
\item The non-Hamiltonian cubic graphs tried all have vector Hamilton cycles,
of dimension in the hundreds. Do they have an explicit description through the
symmetries of the encoding, and which cubic graphs besides the Petersen graph
have none?
\end{enumerate}

\begin{thebibliography}{AGKS00}

\bibitem[Ba95]{Ba95}
  A.I. Barvinok,
  Problems of distance geometry and convex properties of quadratic maps,
  {\em Discrete \& Computational Geometry\/} 13 (1995) 189--202.

\bibitem[Bo99]{Bo99}
  B. Borchers,
  CSDP, a C library for semidefinite programming,
  {\em Optimization Methods and Software\/} 11 (1999) 613--623.

\bibitem[AGKS00]{AGKS00}
  D. Achlioptas, C. Gomes, H. Kautz and B. Selman,
  Generating satisfiable problem instances,
  in: {\em Proceedings of the 17th National Conference on Artificial
  Intelligence (AAAI-00)\/}, 2000, 256--261.

\bibitem[BP03]{BP03}
  S. Busygin and D.V. Pasechnik,
  On $\overline{\chi}(G)-\alpha(G)>0$ gap recognition and
  $\alpha(G)$-upper bounds,
  {\em Electronic Colloquium on Computational Complexity\/}, Report
  TR03-052, 2003 (\texttt{qcp.tex}).

\bibitem[Bu06]{B02}
  S. Busygin,
  A new trust region technique for the maximum weight clique problem,
  {\em Discrete Applied Mathematics\/} 154(15) (2006) 2080--2096.

\bibitem[Bu07]{Bu07}
  S. Busygin,
  The SAT01 framework for NP problems,
  note, 2003--2007 (\texttt{sat01theta.tex} in the SAT01 repository).

\bibitem[CW]{CW}
  S. Busygin,
  A least squares framework for the maximum weight clique problem,
  note, 2007, updated 2010 (\texttt{CliqueWrapper.tex}).

\bibitem[Co84]{Co84}
  C. Colbourn,
  The complexity of completing partial latin squares,
  {\em Discrete Applied Mathematics\/} 8 (1984) 25--30.

\bibitem[CRST06]{CRST06}
  M. Chudnovsky, N. Robertson, P. Seymour and R. Thomas,
  The strong perfect graph theorem,
  {\em Annals of Mathematics\/} 164 (2006) 51--229.

\bibitem[Ga57]{G57}
  D. Gale,
  A theorem on flows in networks,
  {\em Pacific Journal of Mathematics\/} 7 (1957) 1073--1082.

\bibitem[GLS88]{GLS88}
  M. Gr\"otschel, L. Lov\'asz and A. Schrijver,
  {\em Geometric Algorithms and Combinatorial Optimization\/},
  Springer, 1988.

\bibitem[Ho70]{Ho70}
  A.J. Hoffman,
  On eigenvalues and colorings of graphs,
  in: B. Harris, ed., {\em Graph Theory and its Applications\/},
  Academic Press, 1970, 79--91.

\bibitem[JM96]{JM96}
  M.T. Jacobson and P. Matthews,
  Generating uniformly distributed random latin squares,
  {\em Journal of Combinatorial Designs\/} 4 (1996) 405--437.

\bibitem[Ka72]{Ka72}
  R.M. Karp,
  Reducibility among combinatorial problems,
  in: R.E. Miller and J.W. Thatcher, eds., {\em Complexity of Computer
  Computations\/}, Plenum Press, 1972, 85--103.

\bibitem[Kn94]{K94}
  D.E. Knuth,
  The sandwich theorem,
  {\em Electronic Journal of Combinatorics\/} 1 (1994) \#A1.

\bibitem[Lo79]{L79}
  L. Lov\'asz,
  On the Shannon capacity of a graph,
  {\em IEEE Transactions on Information Theory\/} 25 (1979) 1--7.

\bibitem[MS83]{MS83}
  J.J. Mor\'e and D.C. Sorensen,
  Computing a trust region step,
  {\em SIAM Journal on Scientific and Statistical Computing\/} 4 (1983)
  553--572.

\bibitem[N\"O03]{NO03}
  S. Niskanen and P.R.J. \"Osterg\aa rd,
  {\em Cliquer User's Guide, Version 1.0\/},
  Technical Report T48, Communications Laboratory,
  Helsinki University of Technology, 2003.

\bibitem[Ov92]{O92}
  M.L. Overton,
  Large-scale optimization of eigenvalues,
  {\em SIAM Journal on Optimization\/} 2 (1992) 88--120.

\bibitem[Pa98]{P98}
  G. Pataki,
  On the rank of extreme matrices in semidefinite programs and the
  multiplicity of optimal eigenvalues,
  {\em Mathematics of Operations Research\/} 23 (1998) 339--358.

\end{thebibliography}

\end{document}
